% chapters/ch04.tex
% 第~
% \ref{chap:4}~章 $L^p$ 空间
% 源: 原书 PDF 第 98--139 页 (印刷页 89--130)

\chapter{$L^p$ 空间}
\label{chap:4}

设 $(\Omega,\mathcal{M},\mu)$ 表示一个\emph{测度空间}, 即 $\Omega$ 是一个集合, 并且

\begin{enumerate}[label=(\roman*), leftmargin=2.6em, itemsep=0.25em]
	\item $\mathcal{M}$ 是 $\Omega$ 中的一个 \emph{$\sigma$-代数}, 即 $\mathcal{M}$ 是由
	      $\Omega$ 的子集组成的一个族, 使得:
	      \begin{enumerate}[label=(\alph*), leftmargin=2.4em, itemsep=0.15em]
	      	\item $\emptyset\in\mathcal{M}$,
	      	\item $A\in\mathcal{M}\ \Rightarrow\ A^c\in\mathcal{M}$,
	      	\item $\bigcup_{n=1}^\infty A_n\in\mathcal{M}$, 只要 $A_n\in\mathcal{M}\ \forall n$;
	      \end{enumerate}
	\item $\mu$ 是一个\emph{测度}, 即 $\mu:\mathcal{M}\to[0,\infty]$ 满足
	      \begin{enumerate}[label=(\alph*), leftmargin=2.4em, itemsep=0.15em]
	      	\item $\mu(\emptyset)=0$,
	      	\item $\mu\Bigl(\bigcup_{n=1}^\infty A_n\Bigr)=\sum_{n=1}^\infty\mu(A_n)$,
	      	      只要 $(A_n)$ 是 $\mathcal{M}$ 的成员构成的互不相交的可数族.
	      \end{enumerate}
\end{enumerate}

$\mathcal{M}$ 的成员称为\emph{可测集}.有时我们用 $|A|$ 代替 $\mu(A)$.我们还将
假设——尽管这并非本质的——以下条件成立:

\begin{enumerate}[label=(\roman*), leftmargin=2.6em, start=3, itemsep=0.25em]
	\item $\Omega$ 是 \emph{$\sigma$-有限的}, 即存在 $\mathcal{M}$ 中的可数族
	      $(\Omega_n)$, 使得 $\Omega=\bigcup_{n=1}^\infty\Omega_n$ 且
	      $\mu(\Omega_n)<\infty\ \forall n$.
\end{enumerate}

满足 $\mu(E)=0$ 的集合 $E\in\mathcal{M}$ 称为\emph{零测集}.若一个性质在 $\Omega$
上除一个零测集外处处成立, 则说该性质\emph{几乎处处}成立 (或对\emph{几乎所有的}
$x\in\Omega$ 成立), 简记为 a.e..

我们假定读者熟悉可测函数与可积函数 $f:\Omega\to\R$ 的概念; 参见例如
H.~L.~Royden \bref{Royden}{1}, G.~B.~Folland \bref{Folland}{2}, A.~Knapp \bref{Knapp}{1}, D.~L.~Cohn \bref{Cohn}{1}, A.~Friedman \bref{Friedman}{3}, W.~Rudin \bref{Rudin}{2}, P.~Halmos \bref{Halmos}{1}, E.~Hewitt--K.~Stromberg \bref{HewittStromberg}{1}, R.~Wheeden--A.~Zygmund \bref{WheedenZygmund}{1}, J.~Neveu \bref{Neveu}{1}, P.~Malliavin \bref{Malliavin}{1}, A.~J.~Weir \bref{Weir}{1}, A.~Kolmogorov--S.~Fomin \bref{KolmogorovFomin}{1}, I.~Fonseca--G.~Leoni \bref{FonsecaLeoni}{1}.我们用 $L^1(\Omega,\mu)$, 或简记为
$L^1(\Omega)$ (或干脆记作 $L^1$), 表示由 $\Omega$ 到 $\R$ 的可积函数构成的
空间.

我们常把 $\int_\Omega f\,d\mu$ 写成 $\int f$, 并采用记号
\begin{equation*}
\norm{f}_{L^1}=\norm{f}_1=\int_\Omega\abs{f}\,d\mu=\int\abs{f}.
\end{equation*}
按惯例, 我们把几乎处处相等的两个函数视为同一函数.下面回顾一些基本事实.

\section{人人都必须知道的一些积分结果}
\label{sec:4.1}

\begin{keytheorem}[单调收敛定理, Beppo Levi]
\label{thm:4.1}
设 $(f_n)$ 为 $L^1$ 中的函数序列, 满足
\begin{enumerate}[label=(\alph*), leftmargin=2.2em, itemsep=0.1em]
	\item $f_1\le f_2\le\cdots\le f_n\le f_{n+1}\le\cdots$ 在 $\Omega$ 上几乎处处成立,
	\item $\sup_n\int f_n<\infty$.
\end{enumerate}
则 $f_n(x)$ 在 $\Omega$ 上几乎处处收敛于一个有限极限, 记作 $f(x)$; 函数
$f$ 属于 $L^1$, 且 $\norm{f_n-f}_1\to0$.
\end{keytheorem}

\begin{keytheorem}[控制收敛定理, Lebesgue]
\label{thm:4.2}
设 $(f_n)$ 为 $L^1$ 中的函数序列, 满足
\begin{enumerate}[label=(\alph*), leftmargin=2.2em, itemsep=0.1em]
	\item $f_n(x)\to f(x)$ 在 $\Omega$ 上几乎处处成立,
	\item 存在函数 $g\in L^1$, 使得对一切 $n$ 都有 $\abs{f_n(x)}\le g(x)$
	      在 $\Omega$ 上几乎处处成立.
\end{enumerate}
则 $f\in L^1$, 且 $\norm{f_n-f}_1\to0$.
\end{keytheorem}

\begin{lemma}[Fatou 引理]
\label{lem:4.1}
设 $(f_n)$ 为 $L^1$ 中的函数序列, 满足
\begin{enumerate}[label=(\alph*), leftmargin=2.2em, itemsep=0.1em]
	\item 对一切 $n$ 都有 $f_n\ge0$ 几乎处处成立,
	\item $\sup_n\int f_n<\infty$.
\end{enumerate}
对几乎所有的 $x\in\Omega$, 令 $f(x)=\liminf_{n\to\infty}f_n(x)\le+\infty$.则
$f\in L^1$, 且
\begin{equation*}
\int f\le\liminf_{n\to\infty}\int f_n.
\end{equation*}
\end{lemma}

一个基本的例子是 $\Omega=\R^N$, $\mathcal{M}$ 由 Lebesgue 可测集组成,
$\mu$ 为 $\R^N$ 上的 Lebesgue 测度.

\begin{notation}
我们用 $C_c(\R^N)$ 表示 $\R^N$ 上具有紧支集的连续函数构成的空间, 即
\begin{equation*}
C_c(\R^N)=\{f\in C(\R^N);\ f(x)=0\ \ \forall x\in\R^N\setminus K,\
\text{其中}\ K\ \text{紧}\}.
\end{equation*}
\end{notation}

\begin{theorem}[稠密性]
\label{thm:4.3}
空间 $C_c(\R^N)$ 在 $L^1(\R^N)$ 中稠密; 即
\begin{equation*}
\forall f\in L^1(\R^N)\ \ \forall\varepsilon>0\ \ \exists f_1\in C_c(\R^N)\
\text{使得}\ \norm{f-f_1}_1\le\varepsilon.
\end{equation*}
\end{theorem}

设 $(\Omega_1,\mathcal{M}_1,\mu_1)$ 与 $(\Omega_2,\mathcal{M}_2,\mu_2)$ 为两个
$\sigma$-有限的测度空间.按标准方式可以在 Cartesian 积
$\Omega=\Omega_1\times\Omega_2$ 上定义测度空间的结构
$(\Omega,\mathcal{M},\mu)$.

\begin{theorem}[Tonelli]
\label{thm:4.4}
设 $F(x,y):\Omega_1\times\Omega_2\to\R$ 为可测函数, 满足
\begin{equation*}
\text{(a)}\ \ \int_{\Omega_2}\abs{F(x,y)}\,d\mu_2<\infty\ \ \text{对几乎所有的}
\ x\in\Omega_1,
\end{equation*}
以及
\begin{equation*}
\text{(b)}\ \ \int_{\Omega_1}d\mu_1\int_{\Omega_2}\abs{F(x,y)}\,d\mu_2<\infty.
\end{equation*}
则 $F\in L^1(\Omega_1\times\Omega_2)$.
\end{theorem}

\begin{theorem}[Fubini]
\label{thm:4.5}
假设 $F\in L^1(\Omega_1\times\Omega_2)$.则对几乎所有的 $x\in\Omega_1$,
$F(x,y)\in L^1_y(\Omega_2)$, 且 $\int_{\Omega_2}F(x,y)\,d\mu_2\in
L^1_x(\Omega_1)$.同样, 对几乎所有的 $y\in\Omega_2$,
$F(x,y)\in L^1_x(\Omega_1)$, 且 $\int_{\Omega_1}F(x,y)\,d\mu_1\in
L^1_y(\Omega_2)$.此外还有
\begin{equation*}
\int_{\Omega_1}d\mu_1\int_{\Omega_2}F(x,y)\,d\mu_2
=\int_{\Omega_2}d\mu_2\int_{\Omega_1}F(x,y)\,d\mu_1
=\iint_{\Omega_1\times\Omega_2}F(x,y)\,d\mu_1d\mu_2.
\end{equation*}
\end{theorem}

\section{$L^p$ 空间的定义与初等性质}
\label{sec:4.2}

\begin{definition}
设 $p\in\R$, $1<p<\infty$; 令
\begin{equation*}
L^p(\Omega)=\bigl\{f:\Omega\to\R;\ f\ \text{可测且}\ \abs{f}^p\in L^1(\Omega)\bigr\}
\end{equation*}
并置
\begin{equation*}
\norm{f}_{L^p}=\norm{f}_p=\Bigl(\int_\Omega\abs{f(x)}^p\,d\mu\Bigr)^{1/p}.
\end{equation*}
\end{definition}

稍后我们将验证 $\norm{\cdot}_p$ 是一个范数.

\begin{definition}
令
\begin{equation*}
L^\infty(\Omega)=\Bigl\{f:\Omega\to\R\ \Big|\ f\ \text{可测, 且存在常数}\ C\
\text{使得}\ \abs{f(x)}\le C\ \text{在}\ \Omega\ \text{上几乎处处成立}\Bigr\},
\end{equation*}
并置
\begin{equation*}
\norm{f}_{L^\infty}=\norm{f}_\infty=\inf\{C;\ \abs{f(x)}\le C\
\text{在}\ \Omega\ \text{上几乎处处成立}\}.
\end{equation*}
\end{definition}

下面的注记说明 $\norm{\cdot}_\infty$ 是一个范数:

\begin{remark}
\label{rem:4.1}
若 $f\in L^\infty$, 则
\begin{equation*}
\abs{f(x)}\le\norm{f}_\infty\quad\text{在}\ \Omega\ \text{上几乎处处成立}.
\end{equation*}
事实上, 存在序列 $C_n$, 使 $C_n\to\norm{f}_\infty$, 且对每个 $n$ 都有
$\abs{f(x)}\le C_n$ 在 $\Omega$ 上几乎处处成立.于是存在零测集 $E_n$, 使得
$\abs{f(x)}\le C_n$ 对一切 $x\in\Omega\setminus E_n$ 成立.令
$E=\bigcup_{n=1}^\infty E_n$, 则 $|E|=0$, 且
\begin{equation*}
\abs{f(x)}\le C_n\quad\forall n,\ \ \forall x\in\Omega\setminus E;
\end{equation*}
由此推出 $\abs{f(x)}\le\norm{f}_\infty\ \forall x\in\Omega\setminus E$.
\end{remark}

\begin{notation}
设 $1\le p\le\infty$; 我们用 $p'$ 表示\emph{共轭指数},
\begin{equation*}
\boxed{\frac1p+\frac1{p'}=1.}
\end{equation*}
\end{notation}

\index{Hlderbudengshi@Hölder 不等式 (Holder inequality)}
\begin{keytheorem}[Hölder 不等式]
\label{thm:4.6}
假设 $f\in L^p$, $g\in L^{p'}$, 其中 $1\le p\le\infty$.则 $fg\in L^1$, 且
\begin{equation}
\boxed{\int\abs{fg}\le\norm{f}_p\norm{g}_{p'}.}
\label{eq:4.1}
\end{equation}
\end{keytheorem}

\begin{proof}
若 $p=1$ 或 $p=\infty$, 结论显然; 因此设 $1<p<\infty$.我们回顾 Young 不等式:\footnote{有时方便使用
的形式是 $ab\le\varepsilon a^p+C_\varepsilon b^{p'}$, 其中
$C_\varepsilon=\varepsilon^{-1/(p-1)}$.}
\begin{equation}
\boxed{ab\le\frac1pa^p+\frac1{p'}b^{p'}\quad\forall a\ge0,\ \ \forall b\ge0.}
\label{eq:4.2}
\end{equation}
不等式 \eqref{eq:4.2} 是函数 $\log$ 在 $(0,\infty)$ 上凹这一事实的直接推论:
\begin{equation*}
\log\Bigl(\frac1pa^p+\frac1{p'}b^{p'}\Bigr)\ge\frac1p\log a^p+\frac1{p'}\log b^{p'}
=\log ab.
\end{equation*}
我们有
\begin{equation*}
\abs{f(x)g(x)}\le\frac1p\abs{f(x)}^p+\frac1{p'}\abs{g(x)}^{p'}\quad
\text{a.e.}\ x\in\Omega.
\end{equation*}
由此推出 $fg\in L^1$, 且
\begin{equation}
\int\abs{fg}\le\frac1p\norm{f}_p^p+\frac1{p'}\norm{g}_{p'}^{p'}.
\label{eq:4.3}
\end{equation}
在 \eqref{eq:4.3} 中把 $f$ 换成 $\lambda f$ ($\lambda>0$), 得到
\begin{equation}
\int\abs{fg}\le\frac{\lambda^{p-1}}{p}\norm{f}_p^p
+\frac1{\lambda^{p'}p'}\norm{g}_{p'}^{p'}.
\label{eq:4.4}
\end{equation}
选取 $\lambda=\norm{f}_p^{-1}\norm{g}_{p'}^{p'/p}$ (使 \eqref{eq:4.4} 的右端达到
最小), 即得 \eqref{eq:4.1}.
\end{proof}

\begin{remark}
\label{rem:4.2}
记住 Hölder 不等式的下述推广是有益的: 假设 $f_1,f_2,\dots,f_k$ 是满足
\begin{equation*}
f_i\in L^{p_i},\ \ 1\le i\le k,\quad\text{其中}\
\frac1p=\frac1{p_1}+\frac1{p_2}+\cdots+\frac1{p_k}\le1
\end{equation*}
的函数.则乘积 $f=f_1f_2\cdots f_k$ 属于 $L^p$, 且
\begin{equation*}
\norm{f}_p\le\norm{f_1}_{p_1}\norm{f_2}_{p_2}\cdots\norm{f_k}_{p_k}.
\end{equation*}
特别地, 若 $f\in L^p\cap L^q$, 其中 $1\le p\le q\le\infty$, 则对一切满足
$p\le r\le q$ 的 $r$ 都有 $f\in L^r$, 且下述``插值不等式''成立:
\begin{equation*}
\norm{f}_r\le\norm{f}_p^\alpha\norm{f}_q^{1-\alpha},\quad
\text{其中}\ \frac1r=\frac\alpha p+\frac{1-\alpha}q,\ 0\le\alpha\le1;
\end{equation*}
参见习题 4.4.
\end{remark}

\begin{theorem}
\label{thm:4.7}
对任意的 $p$, $1\le p\le\infty$, $L^p$ 是向量空间, 且 $\norm{\cdot}_p$ 是一个范数.
\end{theorem}

\begin{proof}
$p=1$ 与 $p=\infty$ 的情形是清楚的.因此设 $1<p<\infty$, 并设
$f,g\in L^p$.我们有
\begin{equation*}
\abs{f(x)+g(x)}^p\le\bigl(\abs{f(x)}+\abs{g(x)}\bigr)^p
\le2^p\bigl(\abs{f(x)}^p+\abs{g(x)}^p\bigr).
\end{equation*}
从而 $f+g\in L^p$.另一方面,
\begin{equation*}
\norm{f+g}_p^p=\int\abs{f+g}^{p-1}\abs{f+g}
\le\int\abs{f+g}^{p-1}\abs{f}+\int\abs{f+g}^{p-1}\abs{g}.
\end{equation*}
但 $\abs{f+g}^{p-1}\in L^{p'}$, 由 Hölder 不等式得到
\begin{equation*}
\norm{f+g}_p^p\le\norm{f+g}_p^{p-1}\bigl(\norm{f}_p+\norm{g}_p\bigr),
\end{equation*}
即 $\norm{f+g}_p\le\norm{f}_p+\norm{g}_p$.
\end{proof}

\index{FischerRieszdingli@Fischer--Riesz 定理 (Fischer--Riesz theorem)}
\begin{keytheorem}[Fischer--Riesz]
\label{thm:4.8}
对任意的 $p$, $1\le p\le\infty$, $L^p$ 是 Banach 空间.
\end{keytheorem}

\begin{proof}
我们区分 $p=\infty$ 与 $1\le p<\infty$ 两种情形.

\emph{情形 1: $p=\infty$.}设 $(f_n)$ 是 $L^\infty$ 中的 Cauchy 序列.给定整数
$k\ge1$, 存在整数 $N_k$, 使得 $\norm{f_m-f_n}_\infty\le\frac1k$ 对一切
$m,n\ge N_k$ 成立.于是存在零测集 $E_k$, 使得
\begin{equation}
\abs{f_m(x)-f_n(x)}\le\frac1k\quad\forall x\in\Omega\setminus E_k,\ \
\forall m,n\ge N_k.
\label{eq:4.5}
\end{equation}
再令 $E=\bigcup_kE_k$——于是 $E$ 是零测集——我们看出对一切
$x\in\Omega\setminus E$, 序列 $f_n(x)$ 是 Cauchy 序列 (在 $\R$ 中).因此
$f_n(x)\to f(x)$ 对一切 $x\in\Omega\setminus E$ 成立.在 \eqref{eq:4.5} 中令
$m\to\infty$ 取极限, 得到
\begin{equation*}
\abs{f(x)-f_n(x)}\le\frac1k\quad\text{对一切}\ x\in\Omega\setminus E,\ \
\forall n\ge N_k.
\end{equation*}
我们由此断定 $f\in L^\infty$, 且 $\norm{f-f_n}_\infty\le\frac1k$
$\forall n\ge N_k$; 因此 $f_n\to f$ 在 $L^\infty$ 中.

\emph{情形 2: $1\le p<\infty$.}设 $(f_n)$ 是 $L^p$ 中的 Cauchy 序列.为得出结论,
只须证明某个子列在 $L^p$ 中收敛.

我们抽取子列 $(f_{n_k})$, 使得
\begin{equation*}
\norm{f_{n_{k+1}}-f_{n_k}}_p\le\frac1{2^k}\quad\forall k\ge1.
\end{equation*}
[做法如下: 取 $n_1$ 使 $\norm{f_m-f_n}_p\le\frac12$ 对一切 $m,n\ge n_1$ 成立;
再取 $n_2\ge n_1$ 使 $\norm{f_m-f_n}_p\le\frac1{2^2}$ 对一切 $m,n\ge n_2$
成立, 等等.] 我们断言 $f_{n_k}$ 在 $L^p$ 中收敛.为简化记号, 把 $f_{n_k}$
写成 $f_k$, 于是有
\begin{equation}
\norm{f_{k+1}-f_k}_p\le\frac1{2^k}\quad\forall k\ge1.
\label{eq:4.6}
\end{equation}
令
\begin{equation*}
g_n(x)=\sum_{k=1}^n\abs{f_{k+1}(x)-f_k(x)},
\end{equation*}
于是
\begin{equation*}
\norm{g_n}_p\le1.
\end{equation*}
由单调收敛定理, $g_n(x)$ 在 $\Omega$ 上几乎处处趋于一个有限极限, 记作 $g(x)$,
且 $g\in L^p$.另一方面, 对 $m\ge n\ge2$ 有
\begin{equation*}
\abs{f_m(x)-f_n(x)}\le\abs{f_m(x)-f_{m-1}(x)}+\cdots+\abs{f_{n+1}(x)-f_n(x)}
\le g(x)-g_{n-1}(x).
\end{equation*}
由此推出 $f_n(x)$ 在 $\Omega$ 上几乎处处是 Cauchy 序列, 且收敛于一个有限极限,
记作 $f(x)$.在 $\Omega$ 上几乎处处有
\begin{equation}
\abs{f(x)-f_n(x)}\le g(x)\quad\text{对}\ n\ge2,
\label{eq:4.7}
\end{equation}
特别地 $f\in L^p$.最后, 由控制收敛定理推出 $\norm{f_n-f}_p\to0$, 因为
$\abs{f_n(x)-f(x)}^p\to0$ 几乎处处成立, 且
$\abs{f_n-f}^p\le g^p\in L^1$.
\end{proof}

\begin{theorem}
\label{thm:4.9}
设 $(f_n)$ 为 $L^p$ 中的序列, $f\in L^p$ 满足 $\norm{f_n-f}_p\to0$.
则存在子列 $(f_{n_k})$ 与函数 $h\in L^p$, 使得
\begin{enumerate}[label=(\alph*), leftmargin=2.2em, itemsep=0.1em]
	\item $f_{n_k}(x)\to f(x)$ 在 $\Omega$ 上几乎处处成立,
	\item $\abs{f_{n_k}(x)}\le h(x)\ \forall k$, 在 $\Omega$ 上几乎处处成立.
\end{enumerate}
\end{theorem}

\begin{proof}
当 $p=\infty$ 时结论显然.因此设 $1\le p<\infty$.由于 $(f_n)$ 是 Cauchy 序列,
我们可以回到定理  \ref{thm:4.8} 的证明, 考虑一个子列 $(f_{n_k})$——记作
$(f_k)$——它满足 \eqref{eq:4.6}, 并且 $f_k(x)$ 几乎处处收敛于一个极限\footnote{先验地,
应当把 $f$ 与 $f^\star$ 区分开: 按假设 $f_n\to f$ 在 $L^p$ 中, 而另一方面
$f_{n_k}(x)\to f^\star(x)$ 几乎处处成立.} $f^\star(x)$, 其中 $f^\star\in L^p$.
此外, 由 \eqref{eq:4.7} 有 $\abs{f^\star(x)-f_k(x)}\le g(x)\ \forall k$, 在
$\Omega$ 上几乎处处成立, 其中 $g\in L^p$.由控制收敛定理知 $f_k\to f^\star$ 在
$L^p$ 中, 从而 $f=f^\star$ 几乎处处成立.此外还有
$\abs{f_k(x)}\le\abs{f^\star(x)}+g(x)$, 结论随之成立.
\end{proof}

\section{自反性. 可分性. $L^p$ 的对偶}
\label{sec:4.3}

我们分别考虑以下三种情形:
\begin{enumerate}[label=(\Alph*), leftmargin=2.6em, itemsep=0.1em]
	\item $1<p<\infty$,
	\item $p=1$,
	\item $p=\infty$.
\end{enumerate}

\subsection*{A. $1<p<\infty$ 时 $L^p(\Omega)$ 的研究}

这种情形最为``有利'': $L^p$ 是自反的, 可分的, 且 $L^p$ 的对偶是 $L^{p'}$.

\begin{keytheorem}
\label{thm:4.10}
对任意的 $p$, $1<p<\infty$, $L^p$ 是自反的.
\end{keytheorem}

证明分三步.

\runin{第 1 步 (Clarkson 第一不等式).}
设 $2\le p<\infty$.我们断言
\begin{equation}
\Bigl\|\frac{f+g}2\Bigr\|_p^p+\Bigl\|\frac{f-g}2\Bigr\|_p^p
\le\frac12\bigl(\norm{f}_p^p+\norm{g}_p^p\bigr)\quad\forall f,g\in L^p.
\label{eq:4.8}
\end{equation}

\emph{\eqref{eq:4.8} 的证明.}显然只须证明
\begin{equation*}
\Bigl|\frac{a+b}2\Bigr|^p+\Bigl|\frac{a-b}2\Bigr|^p
\le\frac12\bigl(\abs{a}^p+\abs{b}^p\bigr)\quad\forall a,b\in\R.
\end{equation*}
首先注意
\begin{equation*}
\alpha^p+\beta^p\le(\alpha^2+\beta^2)^{p/2}\quad\forall\alpha,\beta\ge0
\end{equation*}
(由齐性可设 $\beta=1$, 并注意函数 $(x^2+1)^{p/2}-x^p-1$ 在 $[0,\infty)$ 上
递增).取 $\alpha=\bigl|\frac{a+b}2\bigr|$, $\beta=\bigl|\frac{a-b}2\bigr|$,
得到
\begin{equation*}
\begin{aligned}
\Bigl|\frac{a+b}2\Bigr|^p+\Bigl|\frac{a-b}2\Bigr|^p
&\le\Bigl[\Bigl|\frac{a+b}2\Bigr|^2+\Bigl|\frac{a-b}2\Bigr|^2\Bigr]^{p/2}
=\Bigl(\frac{a^2}2+\frac{b^2}2\Bigr)^{p/2}\\
&\le\frac12\bigl(\abs{a}^p+\abs{b}^p\bigr)
\end{aligned}
\end{equation*}
(最后一个不等式来自函数 $x\mapsto\abs{x}^{p/2}$ 的凸性, 因为 $p\ge2$).

\runin{第 2 步: $2\le p<\infty$ 时 $L^p$ 一致凸, 从而自反.}
事实上, 设 $\varepsilon>0$, $f,g\in L^p$, 满足 $\norm{f}_p\le1$, $\norm{g}_p\le1$, 且 $\norm{f-g}_p>\varepsilon$.由 \eqref{eq:4.8} 推出
\begin{equation*}
\Bigl\|\frac{f+g}2\Bigr\|_p^p<1-\Bigl(\frac\varepsilon2\Bigr)^p,
\end{equation*}
从而 $\bigl\|\frac{f+g}2\bigr\|_p<1-\delta$, 其中
$\delta=1-\bigl[1-\bigl(\frac\varepsilon2\bigr)^p\bigr]^{1/p}>0$.因此 $L^p$ 一致凸,
由定理  \ref{thm:3.31} 知它自反.

\runin{第 3 步: $1<p\le2$ 时 $L^p$ 自反.}

\begin{proof}
设 $1<p<\infty$.考虑如下定义的算子 $T:L^p\to(L^{p'})^\star$: 固定
$u\in L^p$; 映射 $f\in L^{p'}\mapsto\int uf$ 是 $L^{p'}$ 上的连续线性泛函,
因此它定义了 $(L^{p'})^\star$ 中的一个元素, 记作 $Tu$, 使得
\begin{equation*}
\dual{Tu}{f}=\int uf\quad\forall f\in L^{p'}.
\end{equation*}
我们断言
\begin{equation}
\norm{Tu}_{(L^{p'})^\star}=\norm{u}_p\quad\forall u\in L^p.
\label{eq:4.9}
\end{equation}
事实上, 由 Hölder 不等式有
\begin{equation*}
\abs{\dual{Tu}{f}}\le\norm{u}_p\norm{f}_{p'}\quad\forall f\in L^{p'},
\end{equation*}
因此 $\norm{Tu}_{(L^{p'})^\star}\le\norm{u}_p$.

另一方面, 令
\begin{equation*}
f_0(x)=\abs{u(x)}^{p-2}u(x)
\end{equation*}
(若 $u(x)=0$ 则 $f_0(x)=0$).显然
\begin{equation*}
f_0\in L^{p'},\quad\norm{f_0}_{p'}=\norm{u}_p^{p-1},
\quad\text{且}\quad\dual{Tu}{f_0}=\norm{u}_p^p;
\end{equation*}
于是
\begin{equation}
\norm{Tu}_{(L^{p'})^\star}\ge\frac{\dual{Tu}{f_0}}{\norm{f_0}_{p'}}
=\norm{u}_p.
\label{eq:4.10}
\end{equation}
因此我们证明了 $T$ 是从 $L^p$ 到 $(L^{p'})^\star$ 的等距映射, 这意味着
$T(L^p)$ 是 $(L^{p'})^\star$ 的闭子空间 (因为 $L^p$ 是 Banach 空间).

现在设 $1<p\le2$.由于 $L^{p'}$ 自反 (由第 2 步), 故 $(L^{p'})^\star$ 也自反
(推论  \ref{cor:3.21}).由命题  \ref{prop:3.20} 断定 $T(L^p)$ 自反, 从而 $L^p$
也自反.
\end{proof}

\begin{remark}
\label{rem:4.3}
事实上, $1<p\le2$ 时 $L^p$ 也是一致凸的.这是 Clarkson 第二不等式的推论, 该不等式
对 $1<p\le2$ 成立:
\begin{equation*}
\Bigl\|\frac{f+g}2\Bigr\|_{p'}^p+\Bigl\|\frac{f-g}2\Bigr\|_{p'}^p
\le\Bigl(\frac12\norm{f}_p^p+\frac12\norm{g}_p^p\Bigr)^{1/(p-1)}
\quad\forall f,g\in L^p.
\end{equation*}
这个不等式比 Clarkson 第一不等式更难证明 (参见例如问题 20 或
E.~Hewitt--K.~Stromberg \bref{HewittStromberg}{1}).显然, 它蕴涵 $1<p\le2$ 时
$L^p$ 一致凸; 另一种方法参见 C.~Morawetz \bref{Morawetz}{1} (习题 4.12) 或
J.~Diestel \bref{Diestel}{1}.
\end{remark}

\begin{keytheorem}[Riesz 表示定理]
\label{thm:4.11}
设 $1<p<\infty$, $\varphi\in(L^p)^\star$.则存在唯一的函数 $u\in L^{p'}$,
使得
\begin{equation*}
\dual{\varphi}{f}=\int uf\quad\forall f\in L^p.
\end{equation*}
此外,
\begin{equation*}
\norm{u}_{p'}=\norm{\varphi}_{(L^p)^\star}.
\end{equation*}
\end{keytheorem}

\begin{remark}
\label{rem:4.4}
定理  \ref{thm:4.11} 非常重要.它说明 $1<p<\infty$ 时 $L^p$ 上的每个连续线性泛函
都可以``具体地''表示为一个积分.映射 $\varphi\mapsto u$ 是一个线性满等距, 它使
我们能够把``抽象''空间 $(L^p)^\star$ 与 $L^{p'}$ 等同起来.

在下文中, 我们将系统地作等同
\begin{equation*}
\boxed{(L^p)^\star=L^{p'}.}
\end{equation*}
\end{remark}

\begin{proof}
考虑由 $\dual{Tu}{f}=\int uf\ \forall u\in L^{p'},\ \forall f\in L^p$ 定义的
算子 $T:L^{p'}\to(L^p)^\star$.定理  \ref{thm:4.10} 证明 (第 3 步) 中所用的论证
表明
\begin{equation*}
\norm{Tu}_{(L^p)^\star}=\norm{u}_{p'}\quad\forall u\in L^{p'}.
\end{equation*}
我们断言 $T$ 是满射.事实上, 令 $E=T(L^{p'})$.由于 $E$ 是闭子空间, 只须证明
$E$ 在 $(L^p)^\star$ 中稠密.设 $h\in(L^p)^{\star\star}$ 满足
$\dual{h}{Tu}=0\ \forall u\in L^{p'}$.由于 $L^p$ 自反, $h\in L^p$, 且满足
$\int uh=0\ \forall u\in L^{p'}$.取 $u=\abs{h}^{p-2}h$, 即得 $h=0$.
\end{proof}

\begin{theorem}
\label{thm:4.12}
对任意的 $p$, $1\le p<\infty$, 空间 $C_c(\R^N)$ 在 $L^p(\R^N)$ 中稠密.
\end{theorem}

在证明定理  \ref{thm:4.12} 之前, 先引入一些记号.

\begin{notation}
\emph{截断}运算 $T_n:\R\to\R$ 定义为
\begin{equation*}
T_nr=\begin{cases}
r & \text{若}\ \abs{r}\le n,\\[2pt]
\dfrac{nr}{\abs{r}} & \text{若}\ \abs{r}>n.
\end{cases}
\end{equation*}
给定集合 $E\subset\Omega$, 定义\emph{特征函数}\footnote{不要与第~
\ref{chap:1}~章引入的示性函数
$I_E$ 混淆.} $\chi_E$ 为
\begin{equation*}
\chi_E(x)=\begin{cases}
1 & \text{若}\ x\in E,\\
0 & \text{若}\ x\in\Omega\setminus E.
\end{cases}
\end{equation*}
\end{notation}

\begin{proof}
首先我们断言: 给定 $f\in L^p(\R^N)$ 与 $\varepsilon>0$, 存在函数
$g\in L^\infty(\R^N)$ 与 $\R^N$ 中的紧集 $K$, 使得 $g$ 在 $K$ 之外为零, 且
\begin{equation}
\norm{f-g}_p<\varepsilon.
\label{eq:4.11}
\end{equation}
事实上, 设 $\chi_n$ 为 $B(0,n)$ 的特征函数, 令 $f_n=\chi_nT_nf$.由控制收敛定理
看出 $\norm{f_n-f}_p\to0$, 因此可以取 $n$ 足够大的 $g=f_n$.其次, 给定
$\delta>0$, 由定理  \ref{thm:4.3} 存在 $g_1\in C_c(\R^N)$, 使得
\begin{equation*}
\norm{g-g_1}_1<\delta.
\end{equation*}
我们总可以假设 $\norm{g_1}_\infty\le\norm{g}_\infty$; 否则用 $T_ng_1$ 代替
$g_1$, 其中 $n=\norm{g}_\infty$.最后有
\begin{equation*}
\norm{g-g_1}_p\le\norm{g-g_1}_1^{1/p}\norm{g-g_1}_\infty^{1-(1/p)}
\le\delta^{1/p}\bigl(2\norm{g}_\infty\bigr)^{1-(1/p)}.
\end{equation*}
选取 $\delta>0$ 足够小, 使
\begin{equation*}
\delta^{1/p}\bigl(2\norm{g}_\infty\bigr)^{1-(1/p)}<\varepsilon,
\end{equation*}
即得结论.
\end{proof}

\begin{definition}
若存在 $\mathcal{M}$ 中成员构成的可数族 $(E_n)$, 使得由 $(E_n)$ 生成的
$\sigma$-代数与 $\mathcal{M}$ 重合 (即 $\mathcal{M}$ 是包含一切 $E_n$ 的最小
$\sigma$-代数), 则称测度空间 $\Omega$ 是\emph{可分的}.
\end{definition}

\begin{example}
测度空间 $\Omega=\R^N$ 是可分的.事实上, 对 $(E_n)$ 可以取任意一族可数个开集,
使得 $\R^N$ 中每个开集都能写成这些 $E_n$ 的并.更一般地, 若 $\Omega$ 是可分度量
空间, 而 $\mathcal{M}$ 由 Borel 集组成 (即 $\mathcal{M}$ 是由 $\Omega$ 中的开集
生成的 $\sigma$-代数), 则 $\Omega$ 是可分测度空间.
\end{example}

\begin{theorem}
\label{thm:4.13}
假设 $\Omega$ 是可分测度空间.则对任意的 $p$, $1\le p<\infty$,
$L^p(\Omega)$ 是可分的.
\end{theorem}

我们只考虑 $\Omega=\R^N$ 的情形, 因为一般情形稍嫌麻烦.注意, 作为推论,
对任意可测集 $\Omega\subset\R^N$, $L^p(\Omega)$ 也是可分的.事实上, 存在从
$L^p(\Omega)$ 到 $L^p(\R^N)$ 的典范等距 (在 $\Omega$ 外作零延拓); 因此可以把
$L^p(\Omega)$ 与 $L^p(\R^N)$ 的一个子空间等同起来, 从而 $L^p(\Omega)$ 可分
(由命题  \ref{prop:3.25}).

\begin{proof}[定理 \ref{thm:4.13}, $\Omega=\R^N$ 的情形]
用 $\mathcal{R}$ 表示 $\R^N$ 中形如
$R=\prod_{k=1}^N(a_k,b_k)$ ($a_k,b_k\in\Q$) 的集合构成的可数族.用 $E$ 表示由
函数族 $(\chi_R)_{R\in\mathcal{R}}$ 在 $\Q$ 上生成的向量空间, 即 $E$ 由函数
$\chi_R$ 的, 系数为有理数的有限线性组合组成, 因此 $E$ 是可数的.

我们断言 $E$ 在 $L^p(\R^N)$ 中稠密.事实上, 给定 $f\in L^p(\R^N)$ 与
$\varepsilon>0$, 存在 $f_1\in C_c(\R^N)$, 使得 $\norm{f-f_1}_p<\varepsilon$.
取 $R\in\mathcal{R}$ 为任一包含 $\operatorname{supp}f_1$ ($f_1$ 的支集) 的方体.
给定 $\delta>0$, 容易构造 $f_2\in E$, 使得 $\norm{f_1-f_2}_\infty<\delta$, 且
$f_2$ 在 $R$ 之外为零: 只须把 $R$ 分成 $\mathcal{R}$ 中的小方体, 使 $f_1$ 在
每个小方体上的振幅 (即 $\sup-\inf$) 小于 $\delta$.因此
$\norm{f_1-f_2}_p\le\norm{f_1-f_2}_\infty\abs{R}^{1/p}<\delta\abs{R}^{1/p}$.
只要选取 $\delta>0$ 使 $\delta\abs{R}^{1/p}<\varepsilon$, 便得到
$\norm{f-f_2}_p<2\varepsilon$.
\end{proof}

\subsection*{B. $L^1(\Omega)$ 的研究}

我们先描述 $L^1(\Omega)$ 的对偶空间.

\begin{keytheorem}[Riesz 表示定理]
\label{thm:4.14}
设 $\varphi\in(L^1)^\star$.则存在唯一的函数 $u\in L^\infty$, 使得
\begin{equation*}
\dual{\varphi}{f}=\int uf\quad\forall f\in L^1.
\end{equation*}
此外,
\begin{equation*}
\norm{u}_\infty=\norm{\varphi}_{(L^1)^\star}.
\end{equation*}
\end{keytheorem}

\begin{keyremark}
\label{rem:4.5}
定理  \ref{thm:4.14} 断言 $L^1$ 上的每个连续线性泛函都可以``具体地''表示为一个
积分.映射 $\varphi\mapsto u$ 是一个线性满等距, 它使我们能够把``抽象''空间
$(L^1)^\star$ 与 $L^\infty$ 等同起来.在下文中, 我们将系统地作等同
\begin{equation*}
\boxed{(L^1)^\star=L^\infty.}
\end{equation*}
\end{keyremark}

\begin{proof}
设 $(\Omega_n)$ 为 $\Omega$ 中的可测集序列, 使得
$\Omega=\bigcup_{n=1}^\infty\Omega_n$ 且 $\abs{\Omega_n}<\infty\ \forall n$.令
$\chi_n=\chi_{\Omega_n}$.

$u$ 的唯一性是显然的.事实上, 设 $u\in L^\infty$ 满足
\begin{equation*}
\int uf=0\quad\forall f\in L^1.
\end{equation*}
取 $f=\chi_n\operatorname{sign}u$ (本书通篇采用约定 $\operatorname{sign}0=0$),
即得 $u=0$ 在 $\Omega_n$ 上几乎处处成立, 从而 $u=0$ 在 $\Omega$ 上几乎处处成立.

现在证明 $u$ 的存在性.首先, 我们构造一个函数 $\theta\in L^2(\Omega)$, 使得
\begin{equation*}
\theta(x)\ge\varepsilon_n>0\quad\forall x\in\Omega_n.
\end{equation*}
这样的函数 $\theta$ 的存在性是清楚的.事实上, 令 $\theta$ 在 $\Omega_1$ 上等于
$\alpha_1$, 在 $\Omega_2\setminus\Omega_1$ 上等于 $\alpha_2$, …, 在
$\Omega_n\setminus\Omega_{n-1}$ 上等于 $\alpha_n$, 等等, 并调整常数
$\alpha_n>0$ 使 $\theta\in L^2$.

映射 $f\in L^2(\Omega)\mapsto\dual{\varphi}{\theta f}$ 是 $L^2(\Omega)$ 上的连续
线性泛函.由定理  \ref{thm:4.11} (取 $p=2$) 存在函数 $v\in L^2(\Omega)$, 使得
\begin{equation}
\dual{\varphi}{\theta f}=\int vf\quad\forall f\in L^2(\Omega).
\label{eq:4.12}
\end{equation}
令 $u(x)=v(x)/\theta(x)$.显然 $u$ 有定义, 因为 $\theta>0$ 在 $\Omega$ 上; 此外
$u$ 可测, 且 $u\chi_n\in L^2(\Omega)$.我们断言 $u$ 具有一切所需的性质.有
\begin{equation}
\dual{\varphi}{\chi_ng}=\int u\chi_ng\quad\forall g\in L^\infty(\Omega)\ \ \forall n.
\label{eq:4.13}
\end{equation}
事实上, 只须在 \eqref{eq:4.12} 中取 $f=\chi_ng/\theta$ (注意 $f\in L^2(\Omega)$,
因为 $f$ 在 $\Omega_n$ 上有界且在 $\Omega_n$ 外为零).

其次, 我们断言 $u\in L^\infty(\Omega)$, 且
\begin{equation}
\norm{u}_\infty\le\norm{\varphi}_{(L^1)^\star}.
\label{eq:4.14}
\end{equation}
固定任一常数 $C>\norm{\varphi}_{(L^1)^\star}$, 并令
\begin{equation*}
A=\{x\in\Omega;\ \abs{u(x)}>C\}.
\end{equation*}
我们来验证 $A$ 是零测集.事实上, 在 \eqref{eq:4.13} 中取
$g=\chi_A\operatorname{sign}u$, 得到
\begin{equation*}
\int_{A\cap\Omega_n}\abs{u}\le\norm{\varphi}_{(L^1)^\star}\abs{A\cap\Omega_n},
\end{equation*}
从而
\begin{equation*}
C\abs{A\cap\Omega_n}\le\norm{\varphi}_{(L^1)^\star}\abs{A\cap\Omega_n}.
\end{equation*}
由此推出 $\abs{A\cap\Omega_n}=0\ \forall n$, 从而 $A$ 是零测集.\eqref{eq:4.14}
的证明至此完成.

最后, 我们断言
\begin{equation}
\dual{\varphi}{h}=\int uh\quad\forall h\in L^1(\Omega).
\label{eq:4.15}
\end{equation}
事实上, 只须在 \eqref{eq:4.13} 中取 $g=T_nh$ ($h$ 的截断), 并注意
$\chi_nT_nh\to h$ 在 $L^1(\Omega)$ 中.

为完成定理  \ref{thm:4.14} 的证明, 只剩验证
$\norm{u}_\infty=\norm{\varphi}_{(L^1)^\star}$.由 \eqref{eq:4.15} 有
\begin{equation*}
\abs{\dual{\varphi}{h}}\le\norm{u}_\infty\norm{h}_1\quad\forall h\in L^1(\Omega),
\end{equation*}
因此 $\norm{\varphi}_{(L^1)^\star}\le\norm{u}_\infty$.再结合 \eqref{eq:4.14}
即得结论.
\end{proof}

\begin{keyremark}
\label{rem:4.6}
除 $\Omega$ 由有限个原子组成这一平凡情形 (此时 $L^1(\Omega)$ 是有限维的) 之外,
空间 $L^1(\Omega)$ 绝不是自反的.事实上, 反设 $L^1(\Omega)$ 自反, 并考虑两种
情形:
\begin{enumerate}[label=(\roman*), leftmargin=2.4em, itemsep=0.1em]
	\item $\forall\varepsilon>0\ \exists\omega\subset\Omega$ 可测, 使
	      $0<\mu(\omega)<\varepsilon$.
	\item $\exists\varepsilon>0$, 使得对一切满足 $\mu(\omega)>0$ 的可测集
	      $\omega\subset\Omega$ 都有 $\mu(\omega)\ge\varepsilon$.
\end{enumerate}
在情形 (i) 中, 存在可测集的递减序列 $(\omega_n)$, 使得
$\mu(\omega_n)>0\ \forall n$ 且 $\mu(\omega_n)\to0$ [先任取序列
$(\omega'_k)$, 使 $0<\mu(\omega'_k)<1/2^k$, 再令
$\omega_n=\bigcup_{k=n}^\infty\omega'_k$].令 $\chi_n=\chi_{\omega_n}$, 并定义
$u_n=\chi_n/\norm{\chi_n}_1$.由于 $\norm{u_n}_1=1$, 存在子列——仍记作 $u_n$
——以及某个 $u\in L^1$, 使得 $u_n\weakto u$ 在弱拓扑 $\sigma(L^1,L^\infty)$ 中
(由定理  \ref{thm:3.18}), 即
\begin{equation}
\int u_n\varphi\to\int u\varphi\quad\forall\varphi\in L^\infty.
\label{eq:4.16}
\end{equation}
另一方面, 对固定的 $j$ 与 $n>j$ 有 $\int u_n\chi_j=1$.令 $n\to\infty$ 取极限,
得到 $\int u\chi_j=1\ \forall j$.最后注意 (由控制收敛定理)
$\int u\chi_j\to0$ 当 $j\to\infty$——矛盾.

在情形 (ii) 中, 空间 $\Omega$ 是纯原子的, 且由可数个互不相同的原子 $(a_n)$ 的并
组成 (除非只有有限个原子!).此时 $L^1(\Omega)$ 同构于 $\ell^1$, 故只须证明
$\ell^1$ 不自反.考虑典范基:
\begin{equation*}
e_n=(0,0,\dots,\underset{(n)}1,0,0\dots).
\end{equation*}
假设 $\ell^1$ 自反, 则存在子列 $(e_{n_k})$ 与某个 $x\in\ell^1$, 使得
$e_{n_k}\weakto x$ 在弱拓扑 $\sigma(\ell^1,\ell^\infty)$ 中, 即
\begin{equation*}
\dual{\varphi}{e_{n_k}}\xrightarrow[k\to\infty]{}\dual{\varphi}{x}
\quad\forall\varphi\in\ell^\infty.
\end{equation*}
取
\begin{equation*}
\varphi=\varphi_j=(0,0,\dots,\underset{(j)}1,1,1,\dots),
\end{equation*}
便得到 $\dual{\varphi_j}{x}=1\ \forall j$.另一方面, 当 $j\to\infty$ 时
$\dual{\varphi_j}{x}\to0$ (因为 $x\in\ell^1$)——矛盾.
\end{keyremark}

\subsection*{C. $L^\infty$ 的研究}

我们已经知道 (定理  \ref{thm:4.14}) $L^\infty=(L^1)^\star$.作为对偶空间,
$L^\infty$ 享有若干良好的性质.特别地, 有:

\begin{enumerate}[label=(\roman*), leftmargin=2.6em, itemsep=0.2em]
	\item 闭单位球 $B_{L^\infty}$ 在\wstar{}拓扑 $\sigma(L^\infty,L^1)$ 中紧
	      (由定理  \ref{thm:3.16}).
	\item 若 $\Omega$ 是 $\R^N$ 中的可测子集, $(f_n)$ 是 $L^\infty(\Omega)$ 中的
	      有界序列, 则存在子列 $(f_{n_k})$ 与某个 $f\in L^\infty(\Omega)$, 使得
	      $f_{n_k}\weakstar f$ 在\wstar{}拓扑 $\sigma(L^\infty,L^1)$ 中 (这是推论
	       \ref{cor:3.30} 与定理  \ref{thm:4.13} 的推论).
\end{enumerate}

然而, 除 $\Omega$ 由有限个原子组成这一平凡情形外, $L^\infty(\Omega)$ 不自反;
否则 $L^1(\Omega)$ 将自反 (由推论  \ref{cor:3.21}), 而我们知道 $L^1$ 不自反
(注记 \ref{rem:4.6}).由此推出 $L^\infty$ 的对偶空间 $(L^\infty)^\star$ 包含
$L^1$ (因为 $L^\infty=(L^1)^\star$), 且 $(L^\infty)^\star$ 严格大于 $L^1$.
换言之, 存在 $L^\infty$ 上的连续线性泛函 $\varphi$, 它们不能表示为
\begin{equation*}
\dual{\varphi}{f}=\int uf\quad\forall f\in L^\infty
\end{equation*}
的形式, 其中 $u\in L^1$.

事实上, 我们来描述这样一个泛函的``具体''例子.令
$\varphi_0:C_c(\R^N)\to\R$ 定义为
\begin{equation*}
\varphi_0(f)=f(0)\quad\text{对}\ f\in C_c(\R^N).
\end{equation*}
显然 $\varphi_0$ 是 $C_c(\R^N)$ 上关于 $\norm{\cdot}_\infty$ 范数的连续线性
泛函.由 Hahn--Banach 定理, 可以把 $\varphi_0$ 延拓为 $L^\infty(\R^N)$ 上的
连续线性泛函 $\varphi$, 并且有
\begin{equation}
\dual{\varphi}{f}=f(0)\quad\forall f\in C_c(\R^N).
\label{eq:4.17}
\end{equation}
我们来验证不存在函数 $u\in L^1(\R^N)$ 使得
\begin{equation}
\dual{\varphi}{f}=\int uf\quad\forall f\in L^\infty(\R^N).
\label{eq:4.18}
\end{equation}
反设这样的函数 $u$ 存在.由 \eqref{eq:4.17} 与 \eqref{eq:4.18} 推出
\begin{equation*}
\int uf=0\quad\forall f\in C_c(\R^N)\ \text{且}\ f(0)=0.
\end{equation*}
对 $\Omega=\R^N\setminus\{0\}$ 应用推论  \ref{cor:4.24}, 看出 $u=0$ 在
$\R^N\setminus\{0\}$ 上几乎处处成立, 从而 $u=0$ 在 $\R^N$ 上几乎处处成立.
由 \eqref{eq:4.18} 断定
\begin{equation*}
\dual{\varphi}{f}=0\quad\forall f\in L^\infty(\R^N),
\end{equation*}
这与 \eqref{eq:4.17} 矛盾.

\begin{starremark}
\label{rem:4.7}
$L^\infty$ 的对偶空间与 $L^1$ 并不重合, 但我们仍可以问: $(L^\infty)^\star$ 究竟
是什么样子? 为此, 把 $L^\infty(\Omega;\CC)$ 看作交换 $C^\star$-代数是有益的
(参见例如 W.~Rudin \bref{Rudin}{1}).由 Gelfand 定理,
$L^\infty(\Omega;\CC)$ 同构且等距于某个紧拓扑空间 $K$ 上连续复值函数构成的空间
$C(K;\CC)$ ($K$ 是代数 $L^\infty$ 的谱; 除 $\Omega$ 由有限个原子组成的情形外,
$K$ 不可度量化).因此 $(L^\infty(\Omega;\CC))^\star$ 可以等同于 $K$ 上的复值
Radon 测度构成的空间, 而 $(L^\infty(\Omega;\R))^\star$ 可以等同于 $K$ 上的实值
Radon 测度构成的空间; 更多细节见本章末评注 3, W.~Rudin \bref{Rudin}{1} 以及
K.~Yosida \bref{Yosida}{1} (第 118 页).
\end{starremark}

\begin{remark}
\label{rem:4.8}
除 $\Omega$ 由有限个原子组成的情形外, 空间 $L^\infty(\Omega)$ 不可分.为证明
这一事实, 使用下述结论是方便的.
\end{remark}

\begin{lemma}
\label{lem:4.2}
设 $E$ 为 Banach 空间.假设存在族 $(O_i)_{i\in I}$, 使得
\begin{enumerate}[label=(\roman*), leftmargin=2.4em, itemsep=0.1em]
	\item 对每个 $i\in I$, $O_i$ 是 $E$ 的非空开子集,
	\item 若 $i\ne j$ 则 $O_i\cap O_j=\emptyset$,
	\item $I$ 不可数.
\end{enumerate}
则 $E$ 不可分.
\end{lemma}

\begin{proof}[引理 \ref{lem:4.2}]
反设 $E$ 可分.用 $(u_n)_{n\in\N}$ 表示 $E$ 中一个可数稠密集.对每个
$i\in I$, 集合 $O_i\cap(u_n)_{n\in\N}\ne\emptyset$, 故可以选取 $n(i)$ 使
$u_{n(i)}\in O_i$.映射 $i\mapsto n(i)$ 是单射; 事实上, 若 $n(i)=n(j)$, 则
$u_{n(i)}=u_{n(j)}\in O_i\cap O_j$, 从而 $i=j$.因此 $I$ 可数——矛盾.
\end{proof}

现在证明 $L^\infty(\Omega)$ 不可分.我们断言存在 $\Omega$ 中可测集构成的不可数族
$(\omega_i)_{i\in I}$, 其中各集合互不相同, 即当 $i\ne j$ 时对称差
$\omega_i\triangle\omega_j$ 的测度为正.然后对由
\begin{equation*}
O_i=\{f\in L^\infty(\Omega);\ \norm{f-\chi_{\omega_i}}_\infty<1/2\}
\end{equation*}
定义的族 $(O_i)_{i\in I}$ 应用引理  \ref{lem:4.2} 即得结论 (注意: 若 $\omega$ 与
$\omega'$ 互不相同, 则 $\norm{\chi_\omega-\chi_{\omega'}}_\infty=1$).当 $\Omega$
是 $\R^N$ 中的开集时, 不可数族 $(\omega_i)$ 的存在性是显然的, 因为可以考虑所有
球 $B(x_0,r)$, 其中 $x_0\in\Omega$ 而 $r>0$ 足够小.

当 $\Omega$ 是一般的测度空间时, 把 $\Omega$ 分成它的原子部分 $\Omega_a$ 与非原子
(= 弥散) 部分 $\Omega_d$; 然后区分两种情形:
\begin{enumerate}[label=(\roman*), leftmargin=2.4em, itemsep=0.1em]
	\item $\Omega_d$ 不是零测集.
	\item $\Omega_d$ 是零测集.
\end{enumerate}
在情形 (i) 中, 对每个实数 $t$, $0<t<\mu(\Omega_d)$, 存在可测集 $\omega$ 使
$\mu(\omega)=t$; 参见例如 P.~Halmos \bref{Halmos}{1}, A.~J.~Weir \bref{Weir}{1}
或 J.~Neveu \bref{Neveu}{1}.这样就得到一个不可数的, 互不相同的可测集族.

在情形 (ii) 中, $\Omega$ 由可数个互不相同的原子 $(a_n)$ 的并组成 (除非 $\Omega$
由有限个原子组成).对任意整数集合 $A\subset\N$, 定义
$\omega_A=\bigcup_{n\in A}a_n$.显然 $(\omega_A)$ 是一个不可数的, 互不相同的
可测集族.

下表总结了 $\Omega$ 为 $\R^N$ 的可测子集时空间 $L^p(\Omega)$ 的主要性质:

\begin{center}
\renewcommand{\arraystretch}{1.35}
\begin{tabular}{|c|c|c|c|}
\hline
 & 自反 & 可分 & 对偶空间 \\
\hline
$L^p$, $1<p<\infty$ & 是 & 是 & $L^{p'}$ \\
\hline
$L^1$ & 否 & 是 & $L^\infty$ \\
\hline
$L^\infty$ & 否 & 否 & 严格大于 $L^1$ \\
\hline
\end{tabular}
\end{center}

\section{卷积与正则化}
\label{sec:4.4}

我们先定义函数 $f\in L^1(\R^N)$ 与函数 $g\in L^p(\R^N)$ 的卷积乘积.

\begin{keytheorem}[Young]
\label{thm:4.15}
设 $f\in L^1(\R^N)$, $g\in L^p(\R^N)$, 其中 $1\le p\le\infty$.则对几乎所有的
$x\in\R^N$, 函数 $y\mapsto f(x-y)g(y)$ 在 $\R^N$ 上可积, 且我们定义
\begin{equation*}
\boxed{(f\star g)(x)=\int_{\R^N}f(x-y)g(y)\,dy.}
\end{equation*}
此外 $f\star g\in L^p(\R^N)$, 且
\begin{equation*}
\boxed{\norm{f\star g}_p\le\norm{f}_1\norm{g}_p.}
\end{equation*}
\end{keytheorem}

\begin{proof}
当 $p=\infty$ 时结论显然.考虑两种情形: (i) $p=1$, (ii) $1<p<\infty$.

\emph{情形 (i): $p=1$.}令 $F(x,y)=f(x-y)g(y)$.对几乎所有的 $y\in\R^N$ 有
\begin{equation*}
\int_{\R^N}\abs{F(x,y)}\,dx=\abs{g(y)}\int_{\R^N}\abs{f(x-y)}\,dx
=\abs{g(y)}\norm{f}_1<\infty,
\end{equation*}
并且
\begin{equation*}
\int_{\R^N}dy\int_{\R^N}\abs{F(x,y)}\,dx=\norm{g}_1\norm{f}_1<\infty.
\end{equation*}
由 Tonelli 定理 (定理  \ref{thm:4.4}) 推出 $F\in L^1(\R^N\times\R^N)$.应用
Fubini 定理 (定理  \ref{thm:4.5}), 看出
\begin{equation*}
\int_{\R^N}\abs{F(x,y)}\,dy<\infty\quad\text{对几乎所有的}\ x\in\R^N,
\end{equation*}
并且
\begin{equation*}
\int_{\R^N}dx\int_{\R^N}\abs{F(x,y)}\,dy
=\int_{\R^N}dy\int_{\R^N}\abs{F(x,y)}\,dx=\norm{f}_1\norm{g}_1.
\end{equation*}
这正是 $p=1$ 时定理  \ref{thm:4.15} 的结论.

\emph{情形 (ii): $1<p<\infty$.}由情形 (i) 知, 对几乎所有的固定 $x\in\R^N$,
函数 $y\mapsto\abs{f(x-y)}\abs{g(y)}^p$ 在 $\R^N$ 上可积, 即
\begin{equation*}
\abs{f(x-y)}^{1/p}\abs{g(y)}\in L^p_y(\R^N).
\end{equation*}
由于 $\abs{f(x-y)}^{1/p'}\in L^{p'}_y(\R^N)$, 由 Hölder 不等式推出
\begin{equation*}
\abs{f(x-y)}\abs{g(y)}=\abs{f(x-y)}^{1/p'}\abs{f(x-y)}^{1/p}\abs{g(y)}
\in L^1_y(\R^N),
\end{equation*}
且
\begin{equation*}
\int_{\R^N}\abs{f(x-y)}\abs{g(y)}\,dy
\le\norm{f}_1^{1/p'}\Bigl(\int_{\R^N}\abs{f(x-y)}\abs{g(y)}^p\,dy\Bigr)^{1/p},
\end{equation*}
即
\begin{equation*}
\abs{(f\star g)(x)}^p\le\norm{f}_1^{p/p'}\bigl(\abs{f}\star\abs{g}^p\bigr)(x).
\end{equation*}
由情形 (i) 断定 $f\star g\in L^p(\R^N)$, 且
\begin{equation*}
\norm{f\star g}_p^p\le\norm{f}_1^{p/p'}\norm{f}_1\norm{g}_p^p,
\end{equation*}
即 $\norm{f\star g}_p\le\norm{f}_1\norm{g}_p$.
\end{proof}

\begin{notation}
给定 $\R^N$ 上的函数 $f$, 令 $\check f(x)=f(-x)$.
\end{notation}

\begin{proposition}
\label{prop:4.16}
设 $f\in L^1(\R^N)$, $g\in L^p(\R^N)$, $h\in L^{p'}(\R^N)$.则
\begin{equation*}
\int_{\R^N}(f\star g)h=\int_{\R^N}g(\check f\star h).
\end{equation*}
\end{proposition}

\begin{proof}
函数 $F(x,y)=f(x-y)g(y)h(x)$ 属于 $L^1(\R^N\times\R^N)$, 因为由定理
 \ref{thm:4.15} 与 Hölder 不等式有
\begin{equation*}
\int\abs{h(x)}\,dx\int\abs{f(x-y)}\abs{g(y)}\,dy<\infty.
\end{equation*}
因此
\begin{equation*}
\begin{aligned}
\int(f\star g)(x)h(x)\,dx
&=\int dx\int F(x,y)\,dy=\int dy\int F(x,y)\,dx\\
&=\int g(y)(\check f\star h)(y)\,dy.
\end{aligned}
\end{equation*}
\end{proof}

\runin{支集与卷积.}
函数 $f$ 的支集的概念是标准的: $\operatorname{supp}f$ 是使 $f$ 为零的最大开集的
补集; 换言之, $\operatorname{supp}f$ 是集合 $\{x;\ f(x)\ne0\}$ 的闭包.在处理
等价类 (例如空间 $L^p$) 时, 这个概念并不合适.我们需要一个内蕴的定义, 即当
$f_1=f_2$ 几乎处处成立时, $\operatorname{supp}f_1$ 与 $\operatorname{supp}f_2$
应当相同 (或只相差一个零测集).读者容易承认, 通常的概念对 $\R$ 上的
$f=\chi_\Q$ 没有意义.在下面的命题中我们引入合适的定义.

\begin{proposition}[以及支集的定义]
\label{prop:4.17}
设 $f:\R^N\to\R$ 为任一函数.考虑 $\R^N$ 上满足下述条件的开集构成的族
$(\omega_i)_{i\in I}$: 对每个 $i\in I$, $f=0$ 在 $\omega_i$ 上几乎处处成立.令
$\omega=\bigcup_{i\in I}\omega_i$.则 $f=0$ 在 $\omega$ 上几乎处处成立.

按定义, $\operatorname{supp}f$ 是 $\omega$ 在 $\R^N$ 中的补集.
\end{proposition}

\begin{remark}
\label{rem:4.9}
(a) 假设 $f_1=f_2$ 在 $\R^N$ 上几乎处处成立; 显然
$\operatorname{supp}f_1=\operatorname{supp}f_2$.因此我们可以谈论 $f\in L^p$ 的
$\operatorname{supp}f$——无须指明在等价类中选取了哪个代表元.

(b) 若 $f$ 是 $\R^N$ 上的连续函数, 则容易验证 $\operatorname{supp}f$ 的新定义与
通常的定义一致.
\end{remark}

\begin{proof}[命题 \ref{prop:4.17}]
由于集合 $I$ 不必可数, 所以 $f=0$ 在 $\omega$ 上几乎处处成立这一点并不显然.
但可以如下归到可数情形.存在 $\R^N$ 中开集构成的可数族 $(O_n)$, 使得 $\R^N$ 中
每个开集都是某些 $O_n$ 的并.写 $\omega_i=\bigcup_{n\in A_i}O_n$, 且
$\omega=\bigcup_{n\in B}O_n$, 其中 $B=\bigcup_{i\in I}A_i$.由于对每个
$n\in B$, $f=0$ 在 $O_n$ 上几乎处处成立, 断定 $f=0$ 在 $\omega$ 上几乎处处
成立.
\end{proof}

\begin{keyproposition}
\label{prop:4.18}
设 $f\in L^1(\R^N)$, $g\in L^p(\R^N)$, 其中 $1\le p\le\infty$.则
\begin{equation*}
\boxed{\operatorname{supp}(f\star g)\subset\overline{\operatorname{supp}f
+\operatorname{supp}g}.}
\end{equation*}
\end{keyproposition}

\begin{proof}
固定 $x\in\R^N$, 使函数 $y\mapsto f(x-y)g(y)$ 可积 (见定理  \ref{thm:4.15}).有
\begin{equation*}
(f\star g)(x)=\int f(x-y)g(y)\,dy
=\int_{(x-\operatorname{supp}f)\cap\operatorname{supp}g}f(x-y)g(y)\,dy.
\end{equation*}
若 $x\notin\operatorname{supp}f+\operatorname{supp}g$, 则
$(x-\operatorname{supp}f)\cap\operatorname{supp}g=\emptyset$, 从而
$(f\star g)(x)=0$.于是 $(f\star g)(x)=0$ 在
$(\operatorname{supp}f+\operatorname{supp}g)^c$ 上几乎处处成立.特别地,
\begin{equation*}
(f\star g)(x)=0\quad\text{a.e. 在}\
\operatorname{Int}\bigl[(\operatorname{supp}f+\operatorname{supp}g)^c\bigr]
\ \text{上},
\end{equation*}
因此
\begin{equation*}
\operatorname{supp}(f\star g)\subset\overline{\operatorname{supp}f
+\operatorname{supp}g}.
\end{equation*}
\end{proof}

\begin{keyremark}
\label{rem:4.10}
若 $f$ 与 $g$ 都具有紧支集, 则 $f\star g$ 也具有紧支集.然而, 若两者中只有一个
具有紧支集, 则 $f\star g$ 未必具有紧支集.
\end{keyremark}

\begin{definition}
设 $\Omega\subset\R^N$ 为开集, $1\le p\le\infty$.若对每个紧集 $K\subset\Omega$
都有 $f\chi_K\in L^p(\Omega)$, 则称函数 $f:\Omega\to\R$ 属于
$L^p_{\mathrm{loc}}(\Omega)$.

注意: 若 $f\in L^p_{\mathrm{loc}}(\Omega)$, 则 $f\in L^1_{\mathrm{loc}}(\Omega)$.
\end{definition}

\begin{proposition}
\label{prop:4.19}
设 $f\in C_c(\R^N)$, $g\in L^1_{\mathrm{loc}}(\R^N)$.则 $(f\star g)(x)$ 对每个
$x\in\R^N$ 都有定义, 并且 $(f\star g)\in C(\R^N)$.
\end{proposition}

\begin{proof}
注意对每个 $x\in\R^N$, 函数 $y\mapsto f(x-y)g(y)$ 在 $\R^N$ 上可积, 因此
$(f\star g)(x)$ 对每个 $x\in\R^N$ 都有定义.

设 $x_n\to x$, 并取 $\R^N$ 中固定的紧集 $K$, 使
$(x_n-\operatorname{supp}f)\subset K\ \forall n$.于是
$f(x_n-y)=0\ \forall n,\ \forall y\notin K$.由 $f$ 的一致连续性推出
\begin{equation*}
\abs{f(x_n-y)-f(x-y)}\le\varepsilon_n\chi_K(y)\quad\forall n,\ \forall y\in\R^N,
\end{equation*}
其中 $\varepsilon_n\to0$.于是
\begin{equation*}
\abs{(f\star g)(x_n)-(f\star g)(x)}\le\varepsilon_n\int_K\abs{g(y)}\,dy
\xrightarrow{\ }0.
\end{equation*}
\end{proof}

\begin{notation}
设 $\Omega\subset\R^N$ 为开集.

$C(\Omega)$ 是 $\Omega$ 上连续函数构成的空间.

$C^k(\Omega)$ 是 $\Omega$ 上 $k$ 次连续可微函数构成的空间 ($k\ge1$ 为整数).

$C^\infty(\Omega)=\bigcap_kC^k(\Omega)$.

$C_c(\Omega)$ 是 $\Omega$ 上在 $\Omega$ 内具有紧支集的连续函数构成的空间, 即在某个
紧集 $K\subset\Omega$ 之外为零的函数构成的空间.

$C_c^k(\Omega)=C^k(\Omega)\cap C_c(\Omega)$.

$C_c^\infty(\Omega)=C^\infty(\Omega)\cap C_c(\Omega)$

(有些作者用 $\mathcal{D}(\Omega)$ 或 $C_0^\infty(\Omega)$ 代替
$C_c^\infty(\Omega)$).

若 $f\in C^1(\Omega)$, 其梯度定义为
\begin{equation*}
\nabla f=\Bigl(\frac{\partial f}{\partial x_1},
\frac{\partial f}{\partial x_2},\dots,\frac{\partial f}{\partial x_N}\Bigr).
\end{equation*}
若 $f\in C^k(\Omega)$, 而 $\alpha=(\alpha_1,\alpha_2,\dots,\alpha_N)$ 是长度
$\abs{\alpha}=\alpha_1+\alpha_2+\cdots+\alpha_N$ 小于 $k$ 的多重指标, 则写
\begin{equation*}
D^\alpha f=\frac{\partial^{\alpha_1}}{\partial x_1^{\alpha_1}}
\frac{\partial^{\alpha_2}}{\partial x_2^{\alpha_2}}\cdots
\frac{\partial^{\alpha_N}}{\partial x_N^{\alpha_N}}f.
\end{equation*}
\end{notation}

\begin{keyproposition}
\label{prop:4.20}
设 $f\in C_c^k(\R^N)$ ($k\ge1$), $g\in L^1_{\mathrm{loc}}(\R^N)$.则
$f\star g\in C^k(\R^N)$, 且
\begin{equation*}
\boxed{D^\alpha(f\star g)=(D^\alpha f)\star g\quad\forall\alpha,\
\abs{\alpha}\le k.}
\end{equation*}
特别地, 若 $f\in C_c^\infty(\R^N)$, $g\in L^1_{\mathrm{loc}}(\R^N)$, 则
$f\star g\in C^\infty(\R^N)$.
\end{keyproposition}

\begin{proof}
由归纳法只须考虑 $k=1$ 的情形.给定 $x\in\R^N$, 我们断言 $f\star g$ 在 $x$ 处
可微, 且
\begin{equation*}
\nabla(f\star g)(x)=(\nabla f)\star g(x).
\end{equation*}
设 $h\in\R^N$, $\abs{h}<1$.则对一切 $y\in\R^N$,
\begin{equation*}
\begin{aligned}
&\abs{f(x+h-y)-f(x-y)-h\cdot\nabla f(x-y)}\\
&\quad=\Bigl|\int_0^1\bigl[h\cdot\nabla f(x+sh-y)-h\cdot\nabla f(x-y)\bigr]ds
\Bigr|\le\abs{h}\varepsilon(\abs{h}),
\end{aligned}
\end{equation*}
其中 $\varepsilon(\abs{h})\to0$ 当 $\abs{h}\to0$ (因为 $\nabla f$ 在 $\R^N$ 上
一致连续).

取 $\R^N$ 中固定的紧集 $K$, 使其足够大, 以致
$x+B(0,1)-\operatorname{supp}f\subset K$.有
\begin{equation*}
f(x+h-y)-f(x-y)-h\cdot\nabla f(x-y)=0\quad\forall y\notin K,\
\forall h\in B(0,1),
\end{equation*}
因此
\begin{equation*}
\abs{f(x+h-y)-f(x-y)-h\cdot\nabla f(x-y)}\le\abs{h}\varepsilon(\abs{h})\chi_K(y)
\quad\forall y\in\R^N,\ \forall h\in B(0,1).
\end{equation*}
于是对 $h\in B(0,1)$,
\begin{equation*}
\abs{(f\star g)(x+h)-(f\star g)(x)-h\cdot(\nabla f\star g)(x)}
\le\abs{h}\varepsilon(\abs{h})\int_K\abs{g(y)}\,dy.
\end{equation*}
由此推出 $f\star g$ 在 $x$ 处可微, 且
$\nabla(f\star g)(x)=(\nabla f)\star g(x)$.
\end{proof}

\subsection*{磨光核}

\begin{definition}
\emph{\index{moguanghe@磨光核 (mollifier)}
磨光核序列} $(\rho_n)_{n\ge1}$ 是指 $\R^N$ 上满足下述条件的任一函数序列:
\begin{equation*}
\rho_n\in C_c^\infty(\R^N),\quad
\operatorname{supp}\rho_n\subset\overline{B(0,1/n)},\quad
\int\rho_n=1,\quad\rho_n\ge0\ \text{在}\ \R^N\ \text{上}.
\end{equation*}
在下文中, 我们将系统地使用记号 $(\rho_n)$ 表示一个磨光核序列.
\end{definition}

从一个单一的函数 $\rho\in C_c^\infty(\R^N)$ 出发 (要求
$\operatorname{supp}\rho\subset\overline{B(0,1)}$, $\rho\ge0$ 在 $\R^N$ 上, 且
$\rho$ 不恒为零) 容易生成一个磨光核序列——例如函数
\begin{equation*}
\rho(x)=\begin{cases}
e^{1/(\abs{x}^2-1)} & \text{若}\ \abs{x}<1,\\
0 & \text{若}\ \abs{x}>1.
\end{cases}
\end{equation*}
令 $\rho_n(x)=Cn^N\rho(nx)$, 其中 $C=1/\int\rho$, 便得到磨光核序列.

\begin{proposition}
\label{prop:4.21}
假设 $f\in C(\R^N)$.则 $(\rho_n\star f)\xrightarrow[n\to\infty]{}f$ 在 $\R^N$
的紧集上一致收敛.
\end{proposition}

\begin{proof}\footnote{用卷积作正则化的技巧最初由 Leray 与 Friedrichs 引入.}
设 $K\subset\R^N$ 为固定的紧集.给定 $\varepsilon>0$, 存在 $\delta>0$ (依赖于
$K$ 与 $\varepsilon$), 使得
\begin{equation*}
\abs{f(x-y)-f(x)}<\varepsilon\quad\forall x\in K,\ \forall y\in B(0,\delta).
\end{equation*}
对 $x\in\R^N$ 有
\begin{equation*}
\begin{aligned}
(\rho_n\star f)(x)-f(x)
&=\int\bigl[f(x-y)-f(x)\bigr]\rho_n(y)\,dy\\
&=\int_{B(0,1/n)}\bigl[f(x-y)-f(x)\bigr]\rho_n(y)\,dy.
\end{aligned}
\end{equation*}
当 $n>1/\delta$ 且 $x\in K$ 时得到
\begin{equation*}
\abs{(\rho_n\star f)(x)-f(x)}\le\varepsilon\int\rho_n=\varepsilon.
\end{equation*}
\end{proof}

\begin{keytheorem}
\label{thm:4.22}
假设 $f\in L^p(\R^N)$, $1\le p<\infty$.则
$(\rho_n\star f)\xrightarrow[n\to\infty]{}f$ 在 $L^p(\R^N)$ 中.
\end{keytheorem}

\begin{proof}
给定 $\varepsilon>0$, 固定函数 $f_1\in C_c(\R^N)$, 使
$\norm{f-f_1}_p<\varepsilon$ (见定理  \ref{thm:4.12}).由命题  \ref{prop:4.21} 知
$(\rho_n\star f_1)\to f_1$ 在 $\R^N$ 的每个紧集上一致收敛.另一方面, 由命题
 \ref{prop:4.18} 有
\begin{equation*}
\operatorname{supp}(\rho_n\star f_1)\subset B(0,1/n)+\operatorname{supp}f_1
\subset B(0,1)+\operatorname{supp}f_1,
\end{equation*}
这是一个固定的紧集.由此推出
\begin{equation*}
\norm{(\rho_n\star f_1)-f_1}_p\xrightarrow[n\to\infty]{}0.
\end{equation*}
最后, 写
\begin{equation*}
(\rho_n\star f)-f=\bigl[\rho_n\star(f-f_1)\bigr]
+\bigl[(\rho_n\star f_1)-f_1\bigr]+[f_1-f],
\end{equation*}
于是
\begin{equation*}
\norm{(\rho_n\star f)-f}_p\le2\norm{f-f_1}_p+\norm{(\rho_n\star f_1)-f_1}_p
\end{equation*}
(由定理  \ref{thm:4.15}).我们断定
\begin{equation*}
\limsup_{n\to\infty}\norm{(\rho_n\star f)-f}_p\le2\varepsilon
\quad\forall\varepsilon>0,
\end{equation*}
因此 $\lim_{n\to\infty}\norm{(\rho_n\star f)-f}_p=0$.
\end{proof}

\begin{keycorollary}
\label{cor:4.23}
设 $\Omega\subset\R^N$ 为开集.则对任意的 $1\le p<\infty$,
$C_c^\infty(\Omega)$ 在 $L^p(\Omega)$ 中稠密.
\end{keycorollary}

\begin{proof}
给定 $f\in L^p(\Omega)$, 令
\begin{equation*}
\bar f(x)=\begin{cases}
f(x) & \text{若}\ x\in\Omega,\\
0 & \text{若}\ x\in\R^N\setminus\Omega,
\end{cases}
\end{equation*}
于是 $\bar f\in L^p(\R^N)$.

取 $\R^N$ 中紧集构成的序列 $(K_n)$, 使得
\begin{equation*}
\bigcup_{n=1}^\infty K_n=\Omega
\quad\text{且}\quad\dist(K_n,\Omega^c)\ge2/n\quad\forall n.
\end{equation*}
[例如可以取
$K_n=\{x\in\R^N;\ \abs{x}\le n\ \text{且}\ \dist(x,\Omega^c)\ge2/n\}$.]
令 $g_n=\chi_{K_n}\bar f$, $f_n=\rho_n\star g_n$, 于是
\begin{equation*}
\operatorname{supp}f_n\subset B(0,1/n)+K_n\subset\Omega.
\end{equation*}
由此推出 $f_n\in C_c^\infty(\Omega)$.另一方面有
\begin{equation*}
\begin{aligned}
\norm{f_n-f}_{L^p(\Omega)}=\norm{f_n-\bar f}_{L^p(\R^N)}
&\le\norm{(\rho_n\star g_n)-(\rho_n\star\bar f)}_{L^p(\R^N)}
+\norm{(\rho_n\star\bar f)-\bar f}_{L^p(\R^N)}\\
&\le\norm{g_n-\bar f}_{L^p(\R^N)}
+\norm{(\rho_n\star\bar f)-\bar f}_{L^p(\R^N)}.
\end{aligned}
\end{equation*}
最后注意 $\norm{g_n-\bar f}_{L^p(\R^N)}\to0$ (由控制收敛定理),
而 $\norm{(\rho_n\star\bar f)-\bar f}_{L^p(\R^N)}\to0$ (由定理
 \ref{thm:4.22}).我们断定 $\norm{f_n-f}_{L^p(\Omega)}\to0$.
\end{proof}

\begin{corollary}
\label{cor:4.24}
设 $\Omega\subset\R^N$ 为开集, $u\in L^1_{\mathrm{loc}}(\Omega)$ 满足
\begin{equation*}
\int uf=0\quad\forall f\in C_c^\infty(\Omega).
\end{equation*}
则 $u=0$ 在 $\Omega$ 上几乎处处成立.
\end{corollary}

\begin{proof}
设 $g\in L^\infty(\R^N)$, 其支集 $\operatorname{supp}g$ 是包含在 $\Omega$ 中的
紧集.令 $g_n=\rho_n\star g$, 则当 $n$ 足够大时 $g_n\in C_c^\infty(\Omega)$.
因此有
\begin{equation}
\int ug_n=0\quad\forall n.
\label{eq:4.19}
\end{equation}
由于 $g_n\to g$ 在 $L^1(\R^N)$ 中 (由定理  \ref{thm:4.22}), 存在子列——仍记作
$g_n$——使得 $g_n\to g$ 在 $\R^N$ 上几乎处处成立 (见定理  \ref{thm:4.9}).此外,
$\norm{g_n}_{L^\infty(\R^N)}\le\norm{g}_{L^\infty(\R^N)}$.在 \eqref{eq:4.19}
中取极限 (由控制收敛定理), 得到
\begin{equation}
\int ug=0.
\label{eq:4.20}
\end{equation}

设 $K$ 为包含在 $\Omega$ 中的紧集.取函数 $g$ 为
\begin{equation*}
g=\begin{cases}
\operatorname{sign}u & \text{在}\ K\ \text{上},\\
0 & \text{在}\ \R^N\setminus K\ \text{上}.
\end{cases}
\end{equation*}
由 \eqref{eq:4.20} 推出 $\int_K\abs{u}=0$, 从而 $u=0$ 在 $K$ 上几乎处处成立.
由于这对任一紧集 $K\subset\Omega$ 都成立, 断定 $u=0$ 在 $\Omega$ 上几乎处处成立.
\end{proof}

\section{$L^p$ 中强紧性的判别准则}
\label{sec:4.5}

判断 $L^p(\Omega)$ 中的一族函数是否在 $L^p(\Omega)$ 中 (就强拓扑而言) 具有紧闭包,
是重要的.我们回顾: Ascoli--Arzel\`a 定理回答了 $C(K)$——紧度量空间 $K$ 上取值于
$\R$ 的连续函数空间——中的同一问题.

\begin{keytheorem}[Ascoli--Arzel\`a]
\label{thm:4.25}
设 $K$ 为紧度量空间, $H$ 为 $C(K)$ 的有界子集.假设 $H$ 一致等度连续, 即
\begin{equation}
\forall\varepsilon>0\ \exists\delta>0\ \text{使得}\ d(x_1,x_2)<\delta
\Rightarrow\abs{f(x_1)-f(x_2)}<\varepsilon\quad\forall f\in H.
\label{eq:4.21}
\end{equation}
则 $H$ 在 $C(K)$ 中的闭包紧.
\end{keytheorem}

Ascoli--Arzel\`a 定理的证明可参见例如 W.~Rudin \bref{Rudin}{1}, \bref{Rudin}{2}, A.~Knapp \bref{Knapp}{1}, J.~Dixmier \bref{Dixmier}{1}, A.~Friedman \bref{Friedman}{3}, G.~Choquet \bref{Choquet}{1}, K.~Yosida \bref{Yosida}{1}, H.~L.~Royden \bref{Royden}{1}, J.~R.~Munkres \bref{Munkres}{1}, G.~B.~Folland \bref{Folland}{2} 等.

\begin{notation}[函数的平移]
令 $(\tau_hf)(x)=f(x+h)$, $x\in\R^N$, $h\in\R^N$.
\end{notation}

下面这个定理及其推论是 Ascoli--Arzel\`a 定理的``$L^p$ 版本''.

\begin{keytheorem}[Kolmogorov--M. Riesz--Fr\'echet]
\label{thm:4.26}
设 $\mathcal{F}$ 为 $L^p(\R^N)$ 的有界集, $1\le p<\infty$.假设\footnote{假设
\eqref{eq:4.22} 应与 \eqref{eq:4.21} 比较.它是一个``积分型''等度连续性假设.}
\begin{equation}
\lim_{\abs{h}\to0}\norm{\tau_hf-f}_p=0\quad\text{关于}\ f\in\mathcal{F}\
\text{一致},
\label{eq:4.22}
\end{equation}
即 $\forall\varepsilon>0\ \exists\delta>0$ 使得
$\norm{\tau_hf-f}_p<\varepsilon\ \forall f\in\mathcal{F},\ \forall h\in\R^N$,
$\abs{h}<\delta$.

则对任一测度有限的 $\R^N$ 的可测子集 $\Omega$, $\mathcal{F}|_\Omega$ 在
$L^p(\Omega)$ 中的闭包紧.

[这里 $\mathcal{F}|_\Omega$ 表示 $\mathcal{F}$ 中函数在 $\Omega$ 上的限制.]
\end{keytheorem}

证明分四步.

\runin{第 1 步.}我们断言
\begin{equation}
\norm{(\rho_n\star f)-f}_{L^p(\R^N)}\le\varepsilon\quad\forall f\in\mathcal{F},\
\forall n>1/\delta.
\label{eq:4.23}
\end{equation}
事实上, 由 Hölder 不等式有
\begin{equation*}
\begin{aligned}
\abs{(\rho_n\star f)(x)-f(x)}
&\le\int\abs{f(x-y)-f(x)}\rho_n(y)\,dy\\
&\le\Bigl(\int\abs{f(x-y)-f(x)}^p\rho_n(y)\,dy\Bigr)^{1/p}.
\end{aligned}
\end{equation*}
于是
\begin{equation*}
\begin{aligned}
\int\abs{(\rho_n\star f)(x)-f(x)}^p\,dx
&\le\iint\abs{f(x-y)-f(x)}^p\rho_n(y)\,dx\,dy\\
&=\int_{B(0,1/n)}\rho_n(y)\,dy\int\abs{f(x-y)-f(x)}^p\,dx\le\varepsilon^p,
\end{aligned}
\end{equation*}
只要 $1/n<\delta$.

\runin{第 2 步.}我们断言
\begin{equation}
\norm{\rho_n\star f}_{L^\infty(\R^N)}\le C_n\norm{f}_{L^p(\R^N)}
\quad\forall f\in\mathcal{F}
\label{eq:4.24}
\end{equation}
且
\begin{equation}
\abs{(\rho_n\star f)(x_1)-(\rho_n\star f)(x_2)}\le C_n\norm{f}_p\abs{x_1-x_2}
\quad\forall f\in\mathcal{F},\ \forall x_1,x_2\in\R^N,
\label{eq:4.25}
\end{equation}
其中 $C_n$ 只依赖于 $n$.

不等式 \eqref{eq:4.24} 由 Hölder 不等式得到, 其中 $C_n=\norm{\rho_n}_{p'}$.
另一方面, $\nabla(\rho_n\star f)=(\nabla\rho_n)\star f$, 因此
\begin{equation*}
\norm{\nabla(\rho_n\star f)}_{L^\infty(\R^N)}
\le\norm{\nabla\rho_n}_{L^{p'}(\R^N)}\norm{f}_{L^p(\R^N)}.
\end{equation*}
于是得到 \eqref{eq:4.25}, 其中 $C_n=\norm{\nabla\rho_n}_{L^{p'}(\R^N)}$.

\runin{第 3 步.}给定 $\varepsilon>0$ 与测度有限的 $\Omega\subset\R^N$, 存在
$\Omega$ 的有界可测子集 $\omega$, 使得
\begin{equation}
\norm{f}_{L^p(\Omega\setminus\omega)}<\varepsilon\quad\forall f\in\mathcal{F}.
\label{eq:4.26}
\end{equation}
事实上, 写
\begin{equation*}
\norm{f}_{L^p(\Omega\setminus\omega)}
\le\norm{f-(\rho_n\star f)}_{L^p(\R^N)}
+\norm{\rho_n\star f}_{L^p(\Omega\setminus\omega)}.
\end{equation*}
由 \eqref{eq:4.24}, 只须选取 $\omega$ 使 $\abs{\Omega\setminus\omega}$ 足够小.

\runin{第 4 步: 结论.}由于 $L^p(\Omega)$ 完备, 只须 (参见例如 A.~Knapp
\bref{Knapp}{1} 或 J.~R.~Munkres \bref{Munkres}{1} 第~\ref{sec:7.3}~节) 证明
$\mathcal{F}|_\Omega$ 全有界, 即给定任一 $\varepsilon>0$, 存在
$\mathcal{F}|_\Omega$ 的有限个半径为 $\varepsilon$ 的球覆盖.给定
$\varepsilon>0$, 取有界可测集 $\omega$ 使 \eqref{eq:4.26} 成立, 并取
$n>1/\delta$.族 $H=(\rho_n\star\mathcal{F})|_{\bar\omega}$ 满足
Ascoli--Arzel\`a 定理的全部假设 (由第 2 步).因此 $H$ 在 $C(\bar\omega)$ 中的
闭包紧; 从而 $H$ 在 $L^p(\omega)$ 中的闭包也紧.于是可以用有限个
$L^p(\omega)$ 中半径为 $\varepsilon$ 的球覆盖 $H$, 即
\begin{equation*}
H\subset\bigcup_iB(g_i,\varepsilon)\quad\text{其中}\ g_i\in L^p(\omega).
\end{equation*}
考虑由
\begin{equation*}
\bar g_i=\begin{cases}
g_i & \text{在}\ \omega\ \text{上},\\
0 & \text{在}\ \Omega\setminus\omega\ \text{上}
\end{cases}
\end{equation*}
定义的函数 $\bar g_i:\Omega\to\R$, 以及 $L^p(\Omega)$ 中的球
$B(\bar g_i,3\varepsilon)$.

我们断言这些球覆盖 $\mathcal{F}|_\Omega$.事实上, 给定 $f\in\mathcal{F}$, 存在
某个 $i$ 使
\begin{equation*}
\norm{(\rho_n\star f)-g_i}_{L^p(\omega)}<\varepsilon.
\end{equation*}
由于
\begin{equation*}
\norm{f-\bar g_i}_{L^p(\Omega)}^p=\int_{\Omega\setminus\omega}\abs{f}^p
+\int_\omega\abs{f-g_i}^p,
\end{equation*}
由 \eqref{eq:4.26} 有
\begin{equation*}
\begin{aligned}
\norm{f-\bar g_i}_{L^p(\Omega)}
&\le\varepsilon+\norm{f-g_i}_{L^p(\omega)}\\
&\le\varepsilon+\norm{f-(\rho_n\star f)}_{L^p(\R^N)}
+\norm{(\rho_n\star f)-g_i}_{L^p(\omega)}<3\varepsilon.
\end{aligned}
\end{equation*}
我们断定 $\mathcal{F}|_\Omega$ 在 $L^p(\Omega)$ 中的闭包紧.

\begin{remark}
\label{rem:4.11}
当试图证明 $L^p(\Omega)$ ($\Omega$ 有界) 中的一族函数 $\mathcal{F}$ 在
$L^p(\Omega)$ 中具有紧闭包时, 通常方便的做法是先把函数延拓到整个 $\R^N$ 上, 然后
应用定理  \ref{thm:4.26}, 再考虑它们在 $\Omega$ 上的限制.
\end{remark}

\begin{remark}
\label{rem:4.12}
在定理  \ref{thm:4.26} 的假设下, 一般不能断定 $\mathcal{F}$ 本身在
$L^p(\R^N)$ 中具有紧闭包 (可构造例子, 或参见习题 4.33).还需要一个附加假设;
我们下面来叙述它:
\end{remark}

\begin{corollary}
\label{cor:4.27}
设 $\mathcal{F}$ 为 $L^p(\R^N)$ 的有界集, $1\le p<\infty$.假设
\eqref{eq:4.22} 成立, 并且
\begin{equation}
\begin{cases}
\forall\varepsilon>0\ \exists\Omega\subset\R^N\ \text{有界, 可测, 使得}\\
\norm{f}_{L^p(\R^N\setminus\Omega)}<\varepsilon\quad\forall f\in\mathcal{F}.
\end{cases}
\label{eq:4.27}
\end{equation}
则 $\mathcal{F}$ 在 $L^p(\R^N)$ 中的闭包紧.
\end{corollary}

\begin{proof}
给定 $\varepsilon>0$, 取有界可测集 $\Omega\subset\R^N$ 使 \eqref{eq:4.27} 成立.
由定理  \ref{thm:4.26} 知 $\mathcal{F}|_\Omega$ 在 $L^p(\Omega)$ 中的闭包紧.
于是可以用有限个 $L^p(\Omega)$ 中半径为 $\varepsilon$ 的球覆盖
$\mathcal{F}|_\Omega$, 即
\begin{equation*}
\mathcal{F}|_\Omega\subset\bigcup_iB(g_i,\varepsilon)\quad\text{其中}\
g_i\in L^p(\Omega).
\end{equation*}
令
\begin{equation*}
\bar g_i(x)=\begin{cases}
g_i(x) & \text{在}\ \Omega\ \text{中},\\
0 & \text{在}\ \R^N\setminus\Omega\ \text{上}.
\end{cases}
\end{equation*}
显然 $\mathcal{F}$ 被 $L^p(\R^N)$ 中的球 $B(\bar g_i,2\varepsilon)$ 覆盖.
\end{proof}

\begin{remark}
\label{rem:4.13}
推论  \ref{cor:4.27} 的逆命题也成立 (见习题 4.34).因此我们得到了
$L^p(\R^N)$ 中紧集的完全刻画.
\end{remark}

最后给出定理  \ref{thm:4.26} 的一个有用应用:

\begin{corollary}
\label{cor:4.28}
设 $G$ 为 $L^1(\R^N)$ 中固定的函数, 并令
\begin{equation*}
\mathcal{F}=G\star B,
\end{equation*}
其中 $B$ 是 $L^p(\R^N)$ 中的有界集, $1\le p<\infty$.则对任一测度有限的可测集
$\Omega$, $\mathcal{F}|_\Omega$ 在 $L^p(\Omega)$ 中的闭包紧.
\end{corollary}

\begin{proof}
显然 $\mathcal{F}$ 在 $L^p(\R^N)$ 中有界.另一方面, 若写 $f=G\star u$, 其中
$u\in B$, 则
\begin{equation*}
\norm{\tau_hf-f}_p=\norm{(\tau_hG-G)\star u}_p\le C\norm{\tau_hG-G}_1,
\end{equation*}
再借助下述引理即得结论:
\end{proof}

\begin{lemma}
\label{lem:4.3}
设 $G\in L^q(\R^N)$, $1\le q<\infty$.则
\begin{equation*}
\lim_{h\to0}\norm{\tau_hG-G}_q=0.
\end{equation*}
\end{lemma}

\begin{proof}
给定 $\varepsilon>0$, 由定理  \ref{thm:4.12} 存在函数 $G_1\in C_c(\R^N)$, 使得
$\norm{G-G_1}_q<\varepsilon$.写
\begin{equation*}
\begin{aligned}
\norm{\tau_hG-G}_q
&\le\norm{\tau_hG-\tau_hG_1}_q+\norm{\tau_hG_1-G_1}_q+\norm{G_1-G}_q\\
&\le2\varepsilon+\norm{\tau_hG_1-G_1}_q.
\end{aligned}
\end{equation*}
由于 $\lim_{h\to0}\norm{\tau_hG_1-G_1}_q=0$, 看出
\begin{equation*}
\limsup_{h\to0}\norm{\tau_hG-G}_q\le2\varepsilon\quad\forall\varepsilon>0.
\end{equation*}
\end{proof}

\section*{第 4 章评注}
\addcontentsline{toc}{section}{第 4 章评注}
\markright{第 4 章评注}

\cmtitem{1. Egorov 定理.}
积分论的一些基本结果已在 \S\ref{sec:4.1} 中回顾.有一个未被提及的有用结果如下.

\begin{startheorem}[Egorov]
\label{thm:4.29}
假设 $\Omega$ 是测度有限的测度空间.设 $(f_n)$ 为 $\Omega$ 上的可测函数序列,
使得
\begin{equation*}
f_n(x)\to f(x)\quad\text{a.e. 在}\ \Omega\ \text{上}\quad
(\text{其中}\ \abs{f(x)}<\infty\ \text{a.e.}).
\end{equation*}
则 $\forall\varepsilon>0\ \exists A\subset\Omega$ 可测, 使 $\abs{\Omega\setminus A}
<\varepsilon$ 且 $f_n\to f$ 在 $A$ 上一致收敛.
\end{startheorem}

其证明见\index{Vitalishouliandingli@Vitali 收敛定理 (Vitali convergence theorem)}
习题 4.14, P.~Halmos \bref{Halmos}{1}, G.~B.~Folland \bref{Folland}{2}, E.~Hewitt--K.~Stromberg \bref{HewittStromberg}{1}, R.~Wheeden--A.~Zygmund \bref{WheedenZygmund}{1}, K.~Yosida \bref{Yosida}{1}, A.~Friedman \bref{Friedman}{3} 等.

\cmtitem{2. $L^1$ 中的弱紧集.}
由于 $L^1$ 不自反, $L^1$ 的有界集关于弱拓扑 $\sigma(L^1,L^\infty)$ 并不起重要
作用.下述结果给出了 $L^1$ 中弱紧集的一个有用刻画.

\begin{startheorem}[Dunford--Pettis]
\label{thm:4.30}
设 $\mathcal{F}$ 为 $L^1(\Omega)$ 的有界集.则 $\mathcal{F}$ 在弱拓扑
$\sigma(L^1,L^\infty)$ 中具有紧闭包, 当且仅当 $\mathcal{F}$ 等度可积, 即
\begin{equation*}
\text{(a)}\ \begin{cases}
\forall\varepsilon>0\ \exists\delta>0\ \text{使得}\\
\displaystyle\int_A\abs{f}<\varepsilon\quad
\forall A\subset\Omega\ \text{可测, 且}\ \abs{A}<\delta,\ \forall f\in\mathcal{F},
\end{cases}
\end{equation*}
并且
\begin{equation*}
\text{(b)}\ \begin{cases}
\forall\varepsilon>0\ \exists\omega\subset\Omega\ \text{可测, 且}\
\abs{\omega}<\infty\ \text{使得}\\
\displaystyle\int_{\Omega\setminus\omega}\abs{f}<\varepsilon\quad\forall f\in\mathcal{F}.
\end{cases}
\end{equation*}
\end{startheorem}

定理  \ref{thm:4.30} 的证明与讨论见问题 23, 或 N.~Dunford--J.~T.~Schwartz
\bref{DunfordSchwartz}{1}, B.~Beauzamy \bref{Beauzamy}{1}, J.~Diestel \bref{Diestel}{2}, I.~Fonseca--G.~Leoni \bref{FonsecaLeoni}{1};
关于概率论方面, 还可见 J.~Neveu \bref{Neveu}{1}, C.~Dellacherie--P.~A.~Meyer \bref{DellacherieMeyer}{1}; 另见习题 4.36.

\cmtitem{3. Radon 测度.}
正如刚才所指出的, $L^1$ 的有界集不具有紧性性质.为克服这种紧性的缺失, 有时把
$L^1$ 嵌入一个更大的空间——Radon 测度空间——是非常有用的.

例如, 假设 $\Omega$ 是 $\R^N$ 中带 Lebesgue 测度的有界开集.考虑空间
$E=C(\overline\Omega)$, 配以范数 $\norm{u}=\sup_{x\in\overline\Omega}\abs{u(x)}$.
它的对偶空间记作 $\mathcal{M}(\overline\Omega)$, 称为 $\overline\Omega$ 上的
\emph{Radon 测度}空间.$\mathcal{M}(\overline\Omega)$ 上的\wstar{}拓扑有时称为
``淡''拓扑.

我们把 $L^1(\Omega)$ 与 $\mathcal{M}(\overline\Omega)$ 的一个子空间等同起来.为此
引入如下定义的映射 $L^1(\Omega)\to\mathcal{M}(\overline\Omega)$: 给定
$f\in L^1(\Omega)$, 映射 $u\in C(\overline\Omega)\mapsto\int_\Omega fu\,dx$ 是
$C(\overline\Omega)$ 上的连续线性泛函, 记作 $Tf$, 于是
\begin{equation*}
\dual{Tf}{u}_{E^\star,E}=\int_\Omega fu\,dx\quad\forall u\in E.
\end{equation*}
显然 $T$ 是线性的, 而且 $T$ 是等距, 因为
\begin{equation*}
\norm{Tf}_{\mathcal{M}(\overline\Omega)}
=\sup_{\substack{u\in E\\ \norm{u}\le1}}\int_\Omega fu=\norm{f}_1
\end{equation*}
(见习题 4.26).

利用 $T$ 可以把 $L^1(\Omega)$ 与 $\mathcal{M}(\overline\Omega)$ 的一个子空间
等同起来.由于 $\mathcal{M}(\overline\Omega)$ 是可分空间 $C(\overline\Omega)$ 的
对偶空间, 它在\wstar{}拓扑中具有某些紧性性质.特别地, 若 $(f_n)$ 是
$L^1(\Omega)$ 中的有界序列, 则存在子列 $(f_{n_k})$ 与 Radon 测度 $\mu$, 使得
$f_{n_k}\overset{\star}{\rightharpoonup}\mu$ 在\wstar{}拓扑
$\sigma(E^\star,E)$ 中, 即
\begin{equation*}
\int_\Omega f_{n_k}u\to\dual{\mu}{u}\quad\forall u\in C(\overline\Omega).
\end{equation*}
例如, $L^1$ 中的序列可以关于\wstar{}拓扑收敛于一个 Dirac 测度.Radon 测度的
一些进一步的性质见问题 24.

``测度''这一术语由下述结果提供依据, 它把上述定义与集合论意义下标准的测度概念联系
起来:

\begin{theorem}[Riesz 表示定理]
\label{thm:4.31}
设 $\mu$ 为 $\Omega$ 上的 Radon 测度.则存在 $\Omega$ 上唯一的带号 Borel 测度
$\nu$ (即定义在 $\Omega$ 的 Borel 集上的测度), 使得
\begin{equation*}
\dual{\mu}{u}=\int_\Omega u\,d\nu\quad\forall u\in C(\overline\Omega).
\end{equation*}
\end{theorem}

通常方便的做法是把空间 $E=C(\overline\Omega)$ 换成子空间
\begin{equation*}
E_0=\{f\in C(\overline\Omega);\ f=0\ \text{在}\ \overline\Omega\ \text{的边界上}\}.
\end{equation*}
$E_0$ 的对偶记作 $\mathcal{M}(\Omega)$ (以区别于 $\mathcal{M}(\overline\Omega)$).
Riesz 表示定理仍然成立, 只需附加条件 $\abs{\nu}(\overline\Omega\
\text{的边界})=0$.

关于这一广阔而经典的题材, 参见例如 H.~L.~Royden \bref{Royden}{1}, W.~Rudin \bref{Rudin}{2}, G.~B.~Folland \bref{Folland}{2}, A.~Knapp \bref{Knapp}{1}, P.~Malliavin \bref{Malliavin}{1}, P.~Halmos \bref{Halmos}{1}, I.~Fonseca--G.~Leoni \bref{FonsecaLeoni}{1}.

\cmtitem{4. 向量值函数的 Bochner 积分.}
设 $\Omega$ 为测度空间, $E$ 为 Banach 空间.空间 $L^p(\Omega;E)$ 由定义在
$\Omega$ 上, 取值于 $E$ 的, 在某种适当意义下可测且满足
$\int_\Omega\norm{f(x)}^p\,d\mu<\infty$ 的全部函数 $f$ 组成 (当 $p=\infty$ 时作
通常的修改).\S\ref{sec:4.2} 与 \S\ref{sec:4.3} 中描述的大部分性质在对 $E$ 附加一些假设后仍然成立.
例如, 若 $E$ 自反且 $1<p<\infty$, 则 $L^p(\Omega;E)$ 自反, 其对偶空间为
$L^{p'}(\Omega;E^\star)$.更多细节见 K.~Yosida \bref{Yosida}{1}, D.~L.~Cohn \bref{Cohn}{1}, E.~Hille \bref{Hille}{1}, B.~Beauzamy \bref{Beauzamy}{1}, L.~Schwartz \bref{Schwartz}{3}.
当 $\Omega$ 是 $\R$ 中的区间时, 空间 $L^p(\Omega;E)$ 在演化方程的研究中非常
有用 (见第~
\ref{chap:10}~章).

\cmtitem{5. 插值理论.}
开创插值理论的最引人注目的结果是下述定理.

\begin{theorem}[Schur, M. Riesz, Thorin]
\label{thm:4.32}
假设 $\Omega$ 是测度空间, $\abs{\Omega}<\infty$, 且 $T:L^1(\Omega)\to
L^1(\Omega)$ 是有界线性算子, 其范数为 $M_1=\norm{T}_{\mathcal{L}(L^1,L^1)}$.
此外假设 $T:L^\infty(\Omega)\to L^\infty(\Omega)$ 是有界线性算子, 其范数为
$M_\infty=\norm{T}_{\mathcal{L}(L^\infty,L^\infty)}$.

则对一切 $1<p<\infty$, $T$ 是从 $L^p(\Omega)$ 到 $L^p(\Omega)$ 的有界算子,
且其范数 $M_p$ 满足
\begin{equation*}
M_p\le M_1^{1/p}M_\infty^{1/p'}.
\end{equation*}
\end{theorem}

插值理论最初由 I.~Schur, M.~Riesz, G.~O.~Thorin, J.~Marcinkiewicz 与
A.~Zygmund 发现.许多作者作出了决定性的贡献, 其中包括 J.-L.~Lions, J.~Peetre, A.~P.~Calder\'on, E.~Stein 与 E.~Gagliardo.它已成为调和分析
(参见例如 E.~Stein--G.~Weiss \bref{SteinWeiss}{1}, E.~Stein \bref{Stein}{1}, C.~Sadosky \bref{Sadosky}{1}) 与偏微分方程 (参见例如
J.-L.~Lions--E.~Magenes \bref{LionsMagenes}{1}) 中的有用工具.关于这些问题,
还可见 G.~B.~Folland \bref{Folland}{2}, N.~Dunford--J.~T.~Schwartz
\bref{DunfordSchwartz}{1} (第 1 卷第 520 页), J.~Bergh--J.~L\"ofstr\"om
\bref{BerghLofstrom}{1}, M.~Reed--B.~Simon \bref{ReedSimon}{1} (第 2 卷第 27
页) 以及问题 22.

\cmtitem{6. Young 不等式.}
下面是定理  \ref{thm:4.15} 的一个推广.

\begin{theorem}[Young]
\label{thm:4.33}
假设 $f\in L^p(\R^N)$, $g\in L^q(\R^N)$, 其中 $1\le p\le\infty$,
$1\le q\le\infty$, 且
\begin{equation*}
\frac1r=\frac1p+\frac1q-1\ge0.
\end{equation*}
则 $f\star g\in L^r(\R^N)$, 且 $\norm{f\star g}_r\le\norm{f}_p\norm{g}_q$.
\end{theorem}

其证明参见例如习题 4.30.

\cmtnum{7.}
卷积的概念——它可推广到广义函数 (见 L.~Schwartz \bref{Schwartz}{1} 或
A.~Knapp \bref{Knapp}{2})——在偏微分方程理论中起着基本作用.例如, 方程
$P(D)u=f$ 在 $\R^N$ 中 (其中 $P(D)$ 是任一常系数微分算子) 有形式为
$u=E\star f$ 的解, 其中 $E$ 是 $P(D)$ 的基本解 (Malgrange--Ehrenpreis 定理;
另见第~
\ref{chap:1}~章评注 2b).特别地, 方程 $\Delta u=f$ 在 $\R^3$ 中有形式为
$u=E\star f$ 的解, 其中 $E(x)=-(4\pi\abs{x})^{-1}$.

\section*{第 4 章习题}
\addcontentsline{toc}{section}{第 4 章习题}
\markright{第 4 章习题}

除另有说明外, $\Omega$ 表示一个 $\sigma$-有限测度空间.

\begin{exercise}
\label{exer:4.1}
设 $\alpha>0$, $\beta>0$.令
\begin{equation*}
f(x)=\bigl(1+\abs{x}^\alpha\bigr)^{-1}
\bigl(1+\bigl|\log\abs{x}\bigr|^\beta\bigr)^{-1},\quad x\in\R^N.
\end{equation*}
$f$ 在什么条件下属于 $L^p(\R^N)$?
\end{exercise}

\begin{exercise}
\label{exer:4.2}
假设 $\abs{\Omega}<\infty$, $1\le p\le q\le\infty$.证明 $L^q(\Omega)\subset
L^p(\Omega)$ 且嵌入连续.更确切地说, 证明
\begin{equation*}
\norm{f}_p\le\abs{\Omega}^{\frac1p-\frac1q}\norm{f}_q\quad\forall f\in L^q(\Omega).
\end{equation*}
[提示: 使用 Hölder 不等式.]
\end{exercise}

\begin{exercise}
\label{exer:4.3}
\begin{enumerate}[label=\arabic*., leftmargin=2.2em, itemsep=0.35em]
	\item 设 $f,g\in L^p(\Omega)$, $1\le p\le\infty$.证明
	      \begin{equation*}
	      h(x)=\max\{f(x),g(x)\}\in L^p(\Omega).
	      \end{equation*}
	\item 设 $(f_n)$ 与 $(g_n)$ 为 $L^p(\Omega)$ 中的两个序列, $1\le p\le\infty$,
	      使得 $f_n\to f$ 在 $L^p(\Omega)$ 中, $g_n\to g$ 在 $L^p(\Omega)$ 中.
	      令 $h_n=\max\{f_n,g_n\}$, 证明 $h_n\to h$ 在 $L^p(\Omega)$ 中.
	\item 设 $(f_n)$ 为 $L^p(\Omega)$ 中的序列, $1\le p<\infty$; $(g_n)$ 为
	      $L^\infty(\Omega)$ 中的有界序列.假设 $f_n\to f$ 在 $L^p(\Omega)$ 中,
	      且 $g_n\to g$ 几乎处处.证明 $f_ng_n\to fg$ 在 $L^p(\Omega)$ 中.
\end{enumerate}
\end{exercise}

\begin{exercise}
\label{exer:4.4}
\begin{enumerate}[label=\arabic*., leftmargin=2.2em, itemsep=0.35em]
	\item 设 $f_1,f_2,\dots,f_k$ 是 $k$ 个函数, 满足 $f_i\in L^{p_i}(\Omega)\ \forall i$,
	      其中 $1\le p_i\le\infty$, 且 $\sum_{i=1}^k\frac1{p_i}\le1$.令
	      \begin{equation*}
	      f(x)=\prod_{i=1}^kf_i(x).
	      \end{equation*}
	      证明 $f\in L^p(\Omega)$, 其中 $\frac1p=\sum_{i=1}^k\frac1{p_i}$, 且
	      \begin{equation*}
	      \norm{f}_p\le\prod_{i=1}^k\norm{f_i}_{p_i}.
	      \end{equation*}
	      [提示: 先对 $k=2$ 证明, 再用归纳法.]
	\item 由此推出: 若 $f\in L^p(\Omega)\cap L^q(\Omega)$, 其中
	      $1\le p\le\infty$, $1\le q\le\infty$, 则对 $p$ 与 $q$ 之间的每个 $r$
	      都有 $f\in L^r(\Omega)$.更确切地说, 写
	      \begin{equation*}
	      \frac1r=\frac\alpha p+\frac{1-\alpha}q,\quad\text{其中}\ \alpha\in[0,1],
	      \end{equation*}
	      并证明
	      \begin{equation*}
	      \norm{f}_r\le\norm{f}_p^\alpha\norm{f}_q^{1-\alpha}.
	      \end{equation*}
\end{enumerate}
\end{exercise}

\begin{exercise}
\label{exer:4.5}
设 $1\le p<\infty$, $1\le q\le\infty$.
\begin{enumerate}[label=\arabic*., leftmargin=2.2em, itemsep=0.35em]
	\item 证明 $L^1(\Omega)\cap L^\infty(\Omega)$ 是 $L^p(\Omega)$ 的稠密子集.
	\item 证明集合
	      \begin{equation*}
	      \bigl\{f\in L^p(\Omega)\cap L^q(\Omega);\ \norm{f}_q\le1\bigr\}
	      \end{equation*}
	      在 $L^p(\Omega)$ 中闭.
	\item 设 $(f_n)$ 为 $L^p(\Omega)\cap L^q(\Omega)$ 中的序列, $f\in L^p(\Omega)$.
	      假设 $f_n\to f$ 在 $L^p(\Omega)$ 中, 且 $\norm{f_n}_q\le C$.证明
	      $f\in L^r(\Omega)$, 且对 $p$ 与 $q$ 之间的每个 $r$ ($r\ne q$) 都有
	      $f_n\to f$ 在 $L^r(\Omega)$ 中.
\end{enumerate}
\end{exercise}

\begin{exercise}
\label{exer:4.6}
假设 $\abs{\Omega}<\infty$.
\begin{enumerate}[label=\arabic*., leftmargin=2.2em, itemsep=0.35em]
	\item 设 $f\in L^\infty(\Omega)$.证明 $\lim_{p\to\infty}\norm{f}_p=
	      \norm{f}_\infty$.
	\item 设 $f\in\bigcap_{1\le p<\infty}L^p(\Omega)$, 并假设存在常数 $C$, 使得
	      \begin{equation*}
	      \norm{f}_p\le C\quad\forall1\le p<\infty.
	      \end{equation*}
	      证明 $f\in L^\infty(\Omega)$.
	\item 对 $\Omega=(0,1)$, 构造一个函数 $f\in\bigcap_{1\le p<\infty}L^p(\Omega)$
	      的例子, 使得 $f\notin L^\infty(\Omega)$.
\end{enumerate}
\end{exercise}

\begin{exercise}
\label{exer:4.7}
设 $1\le q\le p\le\infty$.设 $a(x)$ 为 $\Omega$ 上的可测函数.假设对每个函数
$u\in L^p(\Omega)$ 都有 $au\in L^q(\Omega)$.

证明 $a\in L^r(\Omega)$, 其中
\begin{equation*}
r=\begin{cases}
\dfrac{pq}{p-q} & \text{若}\ p<\infty,\\[8pt]
q & \text{若}\ p=\infty.
\end{cases}
\end{equation*}
[提示: 使用闭图像定理.]
\end{exercise}

\begin{exercise}
\label{exer:4.8}
设 $X\subset L^1(\Omega)$ 是 $L^1(\Omega)$ 中的闭向量空间.假设
\begin{equation*}
X\subset\bigcup_{1<q\le\infty}L^q(\Omega).
\end{equation*}
\begin{enumerate}[label=\arabic*., leftmargin=2.2em, itemsep=0.35em]
	\item 证明存在某个 $p>1$, 使得 $X\subset L^p(\Omega)$.
	      [提示: 对每个整数 $n\ge1$ 考虑集合
	      \begin{equation*}
	      X_n=\bigl\{f\in X\cap L^{1+(1/n)}(\Omega);\
	      \norm{f}_{1+(1/n)}\le n\bigr\}.]
	      \end{equation*}
	\item 证明存在常数 $C$, 使得
	      \begin{equation*}
	      \norm{f}_p\le C\norm{f}_1\quad\forall f\in X.
	      \end{equation*}
\end{enumerate}
\end{exercise}

\begin{exercise}[Jensen 不等式]
\label{exer:4.9}
假设 $\abs{\Omega}<\infty$.设 $j:\R\to(-\infty,+\infty]$ 为凸的, 下半连续的
函数, $j\not\equiv+\infty$.设 $f\in L^1(\Omega)$ 满足 $f(x)\in D(j)$ 几乎处处,
且 $j(f)\in L^1(\Omega)$.证明
\begin{equation*}
j\Bigl(\frac1{\abs{\Omega}}\int_\Omega f\Bigr)
\le\frac1{\abs{\Omega}}\int_\Omega j(f).
\end{equation*}
\end{exercise}

\begin{exercise}[凸被积函数]
\label{exer:4.10}
假设 $\abs{\Omega}<\infty$.设 $1\le p<\infty$, $j:\R\to\R$ 为凸的连续函数.
考虑由
\begin{equation*}
J(u)=\begin{cases}
\displaystyle\int_\Omega j(u(x))\,dx & \text{若}\ j(u)\in L^1(\Omega),\\[8pt]
+\infty & \text{若}\ j(u)\notin L^1(\Omega)
\end{cases}
\end{equation*}
定义的函数 $J:L^p(\Omega)\to(-\infty,+\infty]$.
\begin{enumerate}[label=\arabic*., leftmargin=2.2em, itemsep=0.35em]
	\item 证明 $J$ 是凸的.
	\item 证明 $J$ 下半连续.
	      [提示: 先考虑 $j\ge0$ 的情形, 并使用 Fatou 引理.]
	\item 证明共轭函数 $J^\star:L^{p'}(\Omega)\to(-\infty,+\infty]$ 由
	      \begin{equation*}
	      J^\star(f)=\begin{cases}
	      \displaystyle\int_\Omega j^\star(f(x))\,dx & \text{若}\ j^\star(f)\in L^1(\Omega),\\[8pt]
	      +\infty & \text{若}\ j^\star(f)\notin L^1(\Omega)
	      \end{cases}
	      \end{equation*}
	      给出.
	      [提示: 当 $1<p<\infty$ 时考虑
	      $J_n(u)=J(u)+\frac1n\int\abs{u}^p$, 并求出 $J_n^\star$.]
	\item 用 $\partial j$ (相应地, $\partial J$) 表示 $j$ (相应地, $J$) 的次微分
	      (见问题 2).设 $u\in L^p(\Omega)$, $f\in L^{p'}(\Omega)$; 证明
	      \begin{equation*}
	      f\in\partial J(u)\iff f(x)\in\partial j(u(x))\quad\text{a.e. 在}\
	      \Omega\ \text{上}.
	      \end{equation*}
\end{enumerate}
\end{exercise}

\begin{exercise}[空间 $L^\alpha(\Omega)$, $0<\alpha<1$]
\label{exer:4.11}
设 $0<\alpha<1$.令
\begin{equation*}
L^\alpha(\Omega)=\bigl\{u:\Omega\to\R;\ u\ \text{可测且}\
\abs{u}^\alpha\in L^1(\Omega)\bigr\}
\end{equation*}
并置
\begin{equation*}
[u]_\alpha=\Bigl(\int\abs{u}^\alpha\Bigr)^{1/\alpha}.
\end{equation*}
\begin{enumerate}[label=\arabic*., leftmargin=2.2em, itemsep=0.35em]
	\item 验证 $L^\alpha$ 是向量空间, 但 $[\cdot]_\alpha$ 不是范数.更确切地说,
	      证明: 若 $u,v\in L^\alpha(\Omega)$, $u\ge0$ 几乎处处, $v\ge0$ 几乎处处,
	      则
	      \begin{equation*}
	      [u+v]_\alpha\ge[u]_\alpha+[v]_\alpha.
	      \end{equation*}
	\item 证明
	      \begin{equation*}
	      [u+v]_\alpha^\alpha\le[u]_\alpha^\alpha+[v]_\alpha^\alpha
	      \quad\forall u,v\in L^\alpha(\Omega).
	      \end{equation*}
\end{enumerate}
\end{exercise}

\begin{exercise}[$1<p\le2$ 时 $L^p$ 一致凸 (C. Morawetz 的方法)]
\label{exer:4.12}
\begin{enumerate}[label=\arabic*., leftmargin=2.2em, itemsep=0.35em]
	\item 设 $1<p<\infty$.证明存在只依赖于 $p$ 的常数 $C$, 使得
	      \begin{equation*}
	      \abs{a-b}^p\le C\bigl(\abs{a}^p+\abs{b}^p\bigr)^{1-s}
	      \Bigl(\abs{a}^p+\abs{b}^p-2\Bigl|\frac{a+b}2\Bigr|^p\Bigr)^s
	      \quad\forall a,b\in\R,
	      \end{equation*}
	      其中 $s=p/2$.
	\item 由此推出 $1<p\le2$ 时 $L^p(\Omega)$ 一致凸.
	      [提示: 使用第 1 问与 Hölder 不等式.]
\end{enumerate}
\end{exercise}

\begin{exercise}
\label{exer:4.13}
\begin{enumerate}[label=\arabic*., leftmargin=2.2em, itemsep=0.35em]
	\item 验证
	      \begin{equation*}
	      \bigl|\abs{a+b}-\abs{a}-\abs{b}\bigr|\le2\abs{b}\quad\forall a,b\in\R.
	      \end{equation*}
	\item 设 $(f_n)$ 为 $L^1(\Omega)$ 中的序列, 满足
	      \begin{enumerate}[label=(\roman*), leftmargin=2.4em, itemsep=0.1em]
	      	\item $f_n(x)\to f(x)$ 几乎处处,
	      	\item $(f_n)$ 在 $L^1(\Omega)$ 中有界, 即 $\norm{f_n}_1\le M\ \forall n$.
	      \end{enumerate}
	      证明 $f\in L^1(\Omega)$, 且
	      \begin{equation*}
	      \lim_{n\to\infty}\int\bigl\{\abs{f_n}-\abs{f_n-f}\bigr\}=\int\abs{f}.
	      \end{equation*}
	      [提示: 在第 1 问中取 $a=f_n-f$, $b=f$, 并考虑序列
	      $\varphi_n=\bigl|\abs{f_n}-\abs{f_n-f}-\abs{f}\bigr|$.]
	\item 设 $(f_n)$ 为 $L^1(\Omega)$ 中的序列, $f$ 为 $L^1(\Omega)$ 中的函数, 满足
	      \begin{enumerate}[label=(\roman*), leftmargin=2.4em, itemsep=0.1em]
	      	\item $f_n(x)\to f(x)$ 几乎处处,
	      	\item $\norm{f_n}_1\to\norm{f}$.
	      \end{enumerate}
	      证明 $\norm{f_n-f}_1=0$.
\end{enumerate}
\end{exercise}

\begin{exercise}[Egorov 与 Vitali 定理]
\label{exer:4.14}
假设 $\abs{\Omega}<\infty$.设 $(f_n)$ 为可测函数序列, 使得 $f_n\to f$ 几乎处处
(其中 $\abs{f}<\infty$ 几乎处处).
\begin{enumerate}[label=\arabic*., leftmargin=2.2em, itemsep=0.35em]
	\item 设 $\alpha>0$ 固定.证明
	      \begin{equation*}
	      \operatorname{meas}\bigl[\abs{f_n-f}>\alpha\bigr]
	      \xrightarrow[n\to\infty]{}0.
	      \end{equation*}
	\item 更确切地说, 令
	      \begin{equation*}
	      S_n(\alpha)=\bigcup_{k\ge n}\bigl[\abs{f_k-f}>\alpha\bigr].
	      \end{equation*}
	      证明 $\abs{S_n(\alpha)}\xrightarrow[n\to\infty]{}0$.
	\item (Egorov). 证明
	      \begin{equation*}
	      \begin{cases}
	      \forall\delta>0\ \ \exists A\subset\Omega\ \text{可测, 使得}\\
	      \abs{A}<\delta\ \text{且}\ f_n\to f\ \text{在}\ \Omega\setminus A\
	      \text{上一致收敛}.
	      \end{cases}
	      \end{equation*}
	      [提示: 给定整数 $m\ge1$, 借助第 2 问证明存在可测集
	      $\Omega_m\subset\Omega$, 使 $\abs{\Omega_m}<\delta/2^m$, 且存在整数
	      $N_m$ 使得
	      \begin{equation*}
	      \abs{f_k(x)-f(x)}<\frac1m\quad\forall k\ge N_m,\
	      \forall x\in\Omega\setminus\Omega_m.]
	      \end{equation*}
	\item (Vitali). 设 $(f_n)$ 为 $L^p(\Omega)$ 中的序列, $1\le p<\infty$.假设
	      \begin{enumerate}[label=(\roman*), leftmargin=2.4em, itemsep=0.1em]
	      	\item $\forall\varepsilon>0\ \exists\delta>0$ 使得
	      	      $\int_A\abs{f_n}^p<\varepsilon\ \forall n$ 且对一切满足
	      	      $\abs{A}<\delta$ 的可测集 $A\subset\Omega$;
	      	\item $f_n\to f$ 几乎处处.
	      \end{enumerate}
	      证明 $f\in L^p(\Omega)$, 且 $f_n\to f$ 在 $L^p(\Omega)$ 中.
\end{enumerate}
\end{exercise}

\begin{exercise}
\label{exer:4.15}
设 $\Omega=(0,1)$.
\begin{enumerate}[label=\arabic*., leftmargin=2.2em, itemsep=0.35em]
	\item 考虑由 $f_n(x)=ne^{-nx}$ 定义的函数序列 $(f_n)$.证明
	      \begin{enumerate}[label=(\roman*), leftmargin=2.4em, itemsep=0.1em]
	      	\item $f_n\to0$ 几乎处处;
	      	\item $f_n$ 在 $L^1(\Omega)$ 中有界;
	      	\item $f_n\not\to0$ 在 $L^1(\Omega)$ 中强收敛;
	      	\item $f_n\not\weakto0$ 弱收敛于 $\sigma(L^1,L^\infty)$.
	      \end{enumerate}
	      更确切地说, 不存在弱收敛于 $\sigma(L^1,L^\infty)$ 的子列.
	\item 设 $1<p<\infty$, 考虑由 $g_n(x)=n^{1/p}e^{-nx}$ 定义的函数序列
	      $(g_n)$.证明
	      \begin{enumerate}[label=(\roman*), leftmargin=2.4em, itemsep=0.1em]
	      	\item $g_n\to0$ 几乎处处;
	      	\item $(g_n)$ 在 $L^p(\Omega)$ 中有界;
	      	\item $g_n\not\to0$ 在 $L^p(\Omega)$ 中强收敛;
	      	\item $g_n\weakto0$ 弱收敛于 $\sigma(L^p,L^{p'})$.
	      \end{enumerate}
\end{enumerate}
\end{exercise}

\begin{exercise}
\label{exer:4.16}
设 $1<p<\infty$.设 $(f_n)$ 为 $L^p(\Omega)$ 中的序列, 满足
\begin{enumerate}[label=(\roman*), leftmargin=2.4em, itemsep=0.1em]
	\item $f_n$ 在 $L^p(\Omega)$ 中有界,
	\item $f_n\to f$ 在 $\Omega$ 上几乎处处.
\end{enumerate}
\begin{enumerate}[label=\arabic*., leftmargin=2.2em, itemsep=0.35em]
	\item 证明 $f_n\weakto f$ 弱收敛于 $\sigma(L^p,L^{p'})$.
	      [提示: 先证明: 若 $f_n\weakto\tilde f$ 弱收敛于
	      $\sigma(L^p,L^{p'})$ 且 $f_n\to f$ 几乎处处, 则 $f=\tilde f$ 几乎处处
	      (使用习题  \ref{exer:3.4}).]
	\item 若把假设 (ii) 换成
	      \begin{enumerate}[label=(\roman*)$'$, leftmargin=3.0em, itemsep=0.1em]
	      	\item $\norm{f_n-f}_1\to0$,
	      \end{enumerate}
	      结论相同.
	\item 现在假设 (i), (ii) 成立, 且 $\abs{\Omega}<\infty$.证明对每个满足
	      $1\le q<p$ 的 $q$ 都有 $\norm{f_n-f}_q\to0$.
	      [提示: 引入截断函数 $T_kf_n$, 或使用 Egorov 定理.]
\end{enumerate}
\end{exercise}

\begin{exercise}[Brezis--Lieb 引理]
\label{exer:4.17}
设 $1<p<\infty$.
\begin{enumerate}[label=\arabic*., leftmargin=2.2em, itemsep=0.35em]
	\item 证明存在常数 $C$ (依赖于 $p$), 使得
	      \begin{equation*}
	      \bigl|\abs{a+b}^p-\abs{a}^p-\abs{b}^p\bigr|
	      \le C\bigl(\abs{a}^{p-1}\abs{b}+\abs{a}\abs{b}^{p-1}\bigr)
	      \quad\forall a,b\in\R.
	      \end{equation*}
	\item 设 $(f_n)$ 为 $L^p(\Omega)$ 中的有界序列, 使得 $f_n\to f$ 在 $\Omega$ 上
	      几乎处处.证明 $f\in L^p(\Omega)$, 且
	      \begin{equation*}
	      \lim_{n\to\infty}\int_\Omega\bigl\{\abs{f_n}^p-\abs{f_n-f}^p\bigr\}
	      =\int_\Omega\abs{f}^p.
	      \end{equation*}
	      [提示: 在第 1 问中取 $a=f_n-f$, $b=f$.注意由习题  \ref{exer:4.16},
	      $\abs{f_n-f}\weakto0$ 在 $L^p$ 中弱收敛, 而 $\abs{f_n-f}^{p-1}
	      \weakto0$ 在 $L^{p'}$ 中弱收敛.]
	\item 由此推出: 若 $(f_n)$ 为 $L^p(\Omega)$ 中的序列, 满足
	      \begin{enumerate}[label=(\roman*), leftmargin=2.4em, itemsep=0.1em]
	      	\item $f_n(x)\to f(x)$ 几乎处处,
	      	\item $\norm{f_n}_p\to\norm{f}_p$,
	      \end{enumerate}
	      则 $\norm{f_n-f}_p\to0$.
	\item 为第 3 问找出另一种方法.
\end{enumerate}
\end{exercise}

\begin{exercise}[Rademacher 函数]
\label{exer:4.18}
设 $1\le p\le\infty$, $f\in L^p_{\mathrm{loc}}(\R)$.假设 $f$ 是 $T$-周期的,
即 $f(x+T)=f(x)$ 对几乎所有的 $x\in\R$ 成立.令
\begin{equation*}
\bar f=\frac1T\int_0^Tf(t)\,dt.
\end{equation*}
考虑由
\begin{equation*}
u_n(x)=f(nx),\quad x\in(0,1)
\end{equation*}
定义的 $L^p(0,1)$ 中的序列 $(u_n)$.
\begin{enumerate}[label=\arabic*., leftmargin=2.2em, itemsep=0.35em]
	\item 证明 $u_n\weakto\bar f$ 在 $L^p(0,1)$ 中关于拓扑 $\sigma(L^p,L^{p'})$
	      弱收敛.
	\item 求 $\lim_{n\to\infty}\norm{u_n-\bar f}_p$.
	\item 考察下述例子:
	      \begin{enumerate}[label=(\roman*), leftmargin=2.4em, itemsep=0.1em]
	      	\item $u_n(x)=\sin nx$,
	      	\item $u_n(x)=f(nx)$, 其中 $f$ 是 1-周期的, 且
	      	      \begin{equation*}
	      	      f(x)=\begin{cases}
	      	      \alpha & \text{若}\ x\in(0,1/2),\\
	      	      \beta & \text{若}\ x\in(1/2,1).
	      	      \end{cases}
	      	      \end{equation*}
	      \end{enumerate}
	      例 (ii) 中的函数称为 \emph{Rademacher 函数}.
\end{enumerate}
\end{exercise}

\begin{exercise}
\label{exer:4.19}
\begin{enumerate}[label=\arabic*., leftmargin=2.2em, itemsep=0.35em]
	\item 设 $(f_n)$ 为 $L^p(\Omega)$ 中的序列, $1<p<\infty$, $f\in L^p(\Omega)$.
	      假设
	      \begin{enumerate}[label=(\roman*), leftmargin=2.4em, itemsep=0.1em]
	      	\item $f_n\weakto f$ 弱收敛于 $\sigma(L^p,L^{p'})$,
	      	\item $\norm{f_n}_p\to\norm{f}_p$.
	      \end{enumerate}
	      证明 $f_n\to f$ 在 $L^p(\Omega)$ 中强收敛.
	\item 构造 $L^1(0,1)$ 中的序列 $(f_n)$, $f_n\ge0$, 使得:
	      \begin{enumerate}[label=(\roman*), leftmargin=2.4em, itemsep=0.1em]
	      	\item $f_n\weakto f$ 弱收敛于 $\sigma(L^1,L^\infty)$,
	      	\item $\norm{f_n}_1\to\norm{f}_1$,
	      	\item $\norm{f_n-f}_1\not\to0$.
	      \end{enumerate}
	      与习题  \ref{exer:4.13} 的结果以及命题  \ref{prop:3.32} 作比较.
\end{enumerate}
\end{exercise}

\begin{exercise}
\label{exer:4.20}
假设 $\abs{\Omega}<\infty$.设 $1\le p<\infty$, $1\le q<\infty$.

设 $a:\R\to\R$ 为连续函数, 使得
\begin{equation*}
\abs{a(t)}\le C\bigl\{\abs{t}^{p/q}+1\bigr\}\quad\forall t\in\R.
\end{equation*}
考虑由
\begin{equation*}
(Au)(x)=a(u(x)),\quad x\in\Omega
\end{equation*}
定义的 (非线性) 映射 $A:L^p(\Omega)\to L^q(\Omega)$.
\begin{enumerate}[label=\arabic*., leftmargin=2.2em, itemsep=0.35em]
	\item 证明 $A$ 从 $L^p(\Omega)$ 强拓扑到 $L^q(\Omega)$ 强拓扑连续.
	\item 取 $\Omega=(0,1)$, 并假设对每个满足 $u_n\weakto u$ 弱收敛于
	      $\sigma(L^p,L^{p'})$ 的序列 $(u_n)$ 都有 $Au_n\weakto Au$ 弱收敛于
	      $\sigma(L^q,L^{q'})$.证明 $a$ 是仿射函数.
	      [提示: 使用 Rademacher 函数; 见习题  \ref{exer:4.18}.]
\end{enumerate}
\end{exercise}

\begin{exercise}
\label{exer:4.21}
给定函数 $u_0:\R\to\R$, 令 $u_n(x)=u_0(x+n)$.
\begin{enumerate}[label=\arabic*., leftmargin=2.2em, itemsep=0.35em]
	\item 假设 $u_0\in L^p(\R)$, $1<p<\infty$.证明 $u_n\weakto0$ 在 $L^p(\R)$
	      中关于弱拓扑 $\sigma(L^p,L^{p'})$ 弱收敛.
	\item 假设 $u_0\in L^\infty(\R)$, 且 $u_0(x)\to0$ 当 $\abs{x}\to\infty$ 在下述
	      弱意义下成立: 对每个 $\delta>0$, 集合 $[\abs{u_0}>\delta]$ 的测度有限.
	      证明 $u_n\overset{\star}{\rightharpoonup}0$ 在 $L^\infty(\R)$ 中关于
	      \wstar{}拓扑 $\sigma(L^\infty,L^1)$ 弱收敛.
	\item 取 $u_0=\chi_{(0,1)}$.证明不存在子列 $(u_{n_k})$ 在 $L^1(\R)$ 中关于
	      $\sigma(L^1,L^\infty)$ 收敛.
\end{enumerate}
\end{exercise}

\begin{exercise}
\label{exer:4.22}
\begin{enumerate}[label=\arabic*., leftmargin=2.2em, itemsep=0.35em]
	\item 设 $(f_n)$ 为 $L^p(\Omega)$ 中的序列, $1<p\le\infty$, $f\in L^p(\Omega)$.
	      证明下述性质等价:
	      \begin{equation*}
	      \text{(A)}\ \ f_n\weakto f\ \text{在}\ \sigma(L^p,L^{p'})\ \text{中},
	      \end{equation*}
	      \begin{equation*}
	      \text{(B)}\ \begin{cases}
	      \norm{f_n}_p\le C\quad\text{且}\\
	      \displaystyle\int_Ef_n\to\int_Ef\quad\forall E\subset\Omega\ \text{可测, 且}\
	      \abs{E}<\infty.
	      \end{cases}
	      \end{equation*}
	\item 若 $p=1$ 且 $\abs{\Omega}<\infty$, 证明 (A) $\iff$ (B).
	\item 假设 $p=1$ 且 $\abs{\Omega}=\infty$.证明 (A) $\Rightarrow$ (B).
	      构造一个例子说明一般地 (B) $\not\Rightarrow$ (A).
	      [提示: 使用习题  \ref{exer:4.21} 第 3 问.]
	\item 设 $(f_n)$ 为 $L^1(\Omega)$ 中的序列, $f\in L^1(\Omega)$, 且
	      $\abs{\Omega}=\infty$.假设
	      \begin{enumerate}[label=(\alph*), leftmargin=2.4em, itemsep=0.1em]
	      	\item $f_n\ge0\ \forall n$, 且 $f\ge0$ 在 $\Omega$ 上几乎处处,
	      	\item $\int_\Omega f_n\to\int_\Omega f$,
	      	\item $\int_Ef_n\to\int_Ef\ \forall E\subset\Omega$ 可测, 且
	      	      $\abs{E}<\infty$.
	      \end{enumerate}
	      证明 $f_n\weakto f$ 在 $L^1(\Omega)$ 中弱收敛于 $\sigma(L^1,L^\infty)$.
	      [提示: 证明 $\int_Ff_n\to\int_Ff\ \forall F\subset\Omega$ 可测, 且
	      $\abs{F}\le\infty$.]
\end{enumerate}
\end{exercise}

\begin{exercise}
\label{exer:4.23}
设 $f:\Omega\to\R$ 为可测函数, $1\le p\le\infty$.本习题的目的是证明集合
\begin{equation*}
C=\bigl\{u\in L^p(\Omega);\ u\ge f\ \text{a.e.}\bigr\}
\end{equation*}
在 $L^p(\Omega)$ 中关于拓扑 $\sigma(L^p,L^{p'})$ 闭.
\begin{enumerate}[label=\arabic*., leftmargin=2.2em, itemsep=0.35em]
	\item 先假设 $1\le p<\infty$.证明 $C$ 在强 $L^p$ 拓扑中凸且闭.由此推出
	      $C$ 在 $\sigma(L^p,L^{p'})$ 中闭.
	\item 取 $p=\infty$, 证明
	      \begin{equation*}
	      C=\Bigl\{u\in L^\infty(\Omega)\ \Big|\ \int u\phi\ge\int f\phi\
	      \forall\phi\in L^1(\Omega)\ \text{满足}\ f\phi\in L^1(\Omega)\
	      \text{且}\ \phi\ge0\ \text{a.e.}\Bigr\}.
	      \end{equation*}
	      [提示: 先假设 $f\in L^\infty(\Omega)$; 一般情形下引入集合
	      $\omega_n=[\abs{f}<n]$.]
	\item 由此推出: 当 $p=\infty$ 时, $C$ 在 $\sigma(L^\infty,L^1)$ 中闭.
	\item 设 $f_1,f_2\in L^\infty(\Omega)$, $f_1\le f_2$ 几乎处处.证明集合
	      \begin{equation*}
	      C=\bigl\{u\in L^\infty(\Omega);\ f_1\le u\le f_2\ \text{a.e.}\bigr\}
	      \end{equation*}
	      在 $L^\infty(\Omega)$ 中关于拓扑 $\sigma(L^\infty,L^1)$ 紧.
\end{enumerate}
\end{exercise}

\begin{exercise}
\label{exer:4.24}
设 $u\in L^\infty(\R^N)$.设 $(\rho_n)$ 为磨光核序列.设 $(\zeta_n)$ 为
$L^\infty(\R^N)$ 中的序列, 使得
\begin{equation*}
\norm{\zeta_n}_\infty\le1\quad\forall n
\quad\text{且}\quad\zeta_n\to\zeta\ \text{a.e. 在}\ \R^N\ \text{上}.
\end{equation*}
令
\begin{equation*}
v_n=\rho_n\star(\zeta_nu)\quad\text{且}\quad v=\zeta u.
\end{equation*}
\begin{enumerate}[label=\arabic*., leftmargin=2.2em, itemsep=0.35em]
	\item 证明 $v_n\overset{\star}{\rightharpoonup}v$ 在 $L^\infty(\R^N)$ 中关于
	      \wstar{}拓扑 $\sigma(L^\infty,L^1)$ 弱收敛.
	\item 证明对每个球 $B$ 都有 $\int_B\abs{v_n-v}\to0$.
\end{enumerate}
\end{exercise}

\begin{exercise}[$L^\infty(\Omega)$ 中函数的正则化]
\label{exer:4.25}
设 $\Omega\subset\R^N$ 为开集.
\begin{enumerate}[label=\arabic*., leftmargin=2.2em, itemsep=0.35em]
	\item 设 $u\in L^\infty(\Omega)$.证明存在 $C_c^\infty(\Omega)$ 中的序列
	      $(u_n)$, 使得
	      \begin{enumerate}[label=(\alph*), leftmargin=2.4em, itemsep=0.1em]
	      	\item $\norm{u_n}_\infty\le\norm{u}_\infty\ \forall n$,
	      	\item $u_n\to u$ 在 $\Omega$ 上几乎处处,
	      	\item $u_n\overset{\star}{\rightharpoonup}u$ 在 $L^\infty(\Omega)$ 中
	      	      关于\wstar{}拓扑 $\sigma(L^\infty,L^1)$ 弱收敛.
	      \end{enumerate}
	\item 若 $u\ge0$ 在 $\Omega$ 上几乎处处, 证明还可以要求
	      \begin{enumerate}[label=(\alph*), leftmargin=2.4em, start=4, itemsep=0.1em]
	      	\item $u_n\ge0$ 在 $\Omega$ 上 $\forall n$.
	      \end{enumerate}
	\item 由此推出 $C_c^\infty(\Omega)$ 在 $L^\infty(\Omega)$ 中关于拓扑
	      $\sigma(L^\infty,L^1)$ 稠密.
\end{enumerate}
\end{exercise}

\begin{exercise}
\label{exer:4.26}
设 $\Omega\subset\R^N$ 为开集, $f\in L^1_{\mathrm{loc}}(\Omega)$.
\begin{enumerate}[label=\arabic*., leftmargin=2.2em, itemsep=0.35em]
	\item 证明 $f\in L^1(\Omega)$ 当且仅当
	      \begin{equation*}
	      A=\sup\Bigl\{\int f\phi;\ \phi\in C_c(\Omega),\
	      \norm{\phi}_\infty\le1\Bigr\}<\infty.
	      \end{equation*}
	      若 $f\in L^1(\Omega)$, 证明 $A=\norm{f}_1$.
	\item 证明 $f^+\in L^1(\Omega)$ 当且仅当
	      \begin{equation*}
	      B=\sup\Bigl\{\int f\phi;\ \phi\in C_c(\Omega),\
	      \norm{\phi}_\infty\le1\ \text{且}\ \phi\ge0\Bigr\}<\infty.
	      \end{equation*}
	      若 $f^+\in L^1(\Omega)$, 证明 $B=\norm{f^+}_1$.
	\item 把 $C_c(\Omega)$ 换成 $C_c^\infty(\Omega)$ 时, 同样的问题.
	\item 由此推出
	      \begin{equation*}
	      \Bigl\{\int f\phi=0\ \forall\phi\in C_c^\infty(\Omega)\Bigr\}
	      \Longrightarrow\bigl[f=0\ \text{a.e.}\bigr]
	      \end{equation*}
	      以及
	      \begin{equation*}
	      \Bigl\{\int f\phi\ge0\ \forall\phi\in C_c^\infty(\Omega),\
	      \phi\ge0\Bigr\}\Longrightarrow\bigl[f\ge0\ \text{a.e.}\bigr].
	      \end{equation*}
\end{enumerate}
\end{exercise}

\begin{exercise}
\label{exer:4.27}
设 $\Omega\subset\R^N$ 为开集.设 $u,v\in L^1_{\mathrm{loc}}(\Omega)$, 且 $u\ne0$
在一个测度为正的集合上几乎处处成立.假设
\begin{equation*}
\Bigl\{\phi\in C_c^\infty(\Omega)\ \text{且}\ \int u\phi>0\Bigr\}
\Longrightarrow\Bigl\{\int v\phi\ge0\Bigr\}.
\end{equation*}
证明存在常数 $\lambda\ge0$, 使得 $v=\lambda u$.
\end{exercise}

\begin{exercise}
\label{exer:4.28}
设 $\rho\in L^1(\R^N)$, $\int\rho=1$.令 $\rho_n(x)=n^N\rho(nx)$.设
$f\in L^p(\R^N)$, $1\le p<\infty$.证明 $\rho_n\star f\to f$ 在 $L^p(\R^N)$ 中.
\end{exercise}

\begin{exercise}
\label{exer:4.29}
设 $K\subset\R^N$ 为紧子集.证明存在 $C_c^\infty(\R^N)$ 中的函数序列 $(u_n)$,
使得
\begin{enumerate}[label=(\alph*), leftmargin=2.4em, itemsep=0.15em]
	\item $0\le u_n\le1$ 在 $\R^N$ 上,
	\item $u_n=1$ 在 $K$ 上,
	\item $\operatorname{supp}u_n\subset K+B(0,1/n)$,
	\item $\abs{D^\alpha u_n(x)}\le C_\alpha n^{\abs{\alpha}}\ \forall x\in\R^N$,
	      对每个多重指标 $\alpha$ (其中 $C_\alpha$ 只依赖于 $\alpha$ 而不依赖于
	      $n$).
\end{enumerate}
[提示: 设 $\chi_n$ 为 $K+B(0,1/2n)$ 的特征函数; 取
$u_n=\rho_{2n}\star\chi_n$.]
\end{exercise}

\begin{exercise}[Young 不等式]
\label{exer:4.30}
设 $1\le p\le\infty$, $1\le q\le\infty$, 满足 $\frac1p+\frac1q\ge1$.令
$\frac1r=\frac1p+\frac1q-1$, 于是 $1\le r\le\infty$.设
$f\in L^p(\R^N)$, $g\in L^q(\R^N)$.
\begin{enumerate}[label=\arabic*., leftmargin=2.2em, itemsep=0.35em]
	\item 证明对几乎所有的 $x\in\R^N$, 函数 $y\mapsto f(x-y)g(y)$ 在 $\R^N$ 上
	      可积.
	      [提示: 令 $\alpha=p/q'$, $\beta=q/p'$, 并写
	      \begin{equation*}
	      \abs{f(x-y)g(y)}=\abs{f(x-y)}^\alpha\abs{g(y)}^\beta
	      \bigl[\abs{f(x-y)}^{1-\alpha}\abs{g(y)}^{1-\beta}\bigr].]
	      \end{equation*}
	\item 令
	      \begin{equation*}
	      (f\star g)(x)=\int_{\R^N}f(x-y)g(y)\,dy.
	      \end{equation*}
	      证明 $f\star g\in L^r(\R^N)$, 且 $\norm{f\star g}_r\le\norm{f}_p
	      \norm{g}_q$.
	\item 这里假设 $\frac1p+\frac1q=1$.证明
	      \begin{equation*}
	      f\star g\in C(\R^N)\cap L^\infty(\R^N),
	      \end{equation*}
	      并且, 若 $1<p<\infty$, 则当 $\abs{x}\to\infty$ 时
	      $(f\star g)(x)\to0$.
\end{enumerate}
\end{exercise}

\begin{exercise}
\label{exer:4.31}
设 $f\in L^p(\R^N)$, $1\le p<\infty$.对每个 $r>0$ 令
\begin{equation*}
f_r(x)=\frac1{\abs{B(x,r)}}\int_{B(x,r)}f(y)\,dy,\quad x\in\R^N.
\end{equation*}
\begin{enumerate}[label=\arabic*., leftmargin=2.2em, itemsep=0.35em]
	\item 证明 $f_r\in L^p(\R^N)\cap C(\R^N)$, 且当 $\abs{x}\to\infty$ 时
	      $f_r(x)\to0$ ($r$ 固定).
	\item 证明当 $r\to0$ 时 $f_r\to f$ 在 $L^p(\R^N)$ 中.
	      [提示: 把 $f_r$ 写成 $\phi_r\star f$ 的形式, 其中 $\phi_r$ 取法适当.]
\end{enumerate}
\end{exercise}

\begin{exercise}
\label{exer:4.32}
\begin{enumerate}[label=\arabic*., leftmargin=2.2em, itemsep=0.35em]
	\item 设 $f,g\in L^1(\R^N)$, $h\in L^p(\R^N)$, $1\le p\le\infty$.证明
	      $f\star g=g\star f$ 且 $(f\star g)\star h=f\star(g\star h)$.
	\item 设 $f\in L^1(\R^N)$.假设 $f\star\phi=0\ \forall\phi\in
	      C_c^\infty(\R^N)$.证明 $f=0$ 在 $\R^N$ 上几乎处处成立.对
	      $f\in L^1_{\mathrm{loc}}(\R^N)$ 也有同样的问题.
	\item 设 $a\in L^1(\R^N)$ 为固定函数.考虑由
	      \begin{equation*}
	      T_a(u)=a\star u
	      \end{equation*}
	      定义的算子 $T_a:L^2(\R^N)\to L^2(\R^N)$.验证 $T_a$ 有界, 且
	      $\norm{T_a}_{\mathcal{L}(L^2)}\le\norm{a}_{L^1(\R^N)}$.计算
	      $T_a\circ T_b$, 并证明 $T_a\circ T_b=T_b\circ T_a\
	      \forall a,b\in L^1(\R^N)$.确定 $(T_a)^\star$, $T_a\circ(T_a)^\star$
	      与 $(T_a)^\star\circ T_a$.在关于 $a$ 的什么条件下有
	      $(T_a)^\star=T_a$?
\end{enumerate}
\end{exercise}

\begin{exercise}
\label{exer:4.33}
固定函数 $\phi\in C_c(\R)$, $\phi\not\equiv0$, 并考虑函数族
\begin{equation*}
\mathcal{F}=\bigcup_{n=1}^\infty\{\phi_n\},
\end{equation*}
其中 $\phi_n(x)=\phi(x+n)$, $x\in\R$.
\begin{enumerate}[label=\arabic*., leftmargin=2.2em, itemsep=0.35em]
	\item 假设 $1\le p<\infty$.证明 $\forall\varepsilon>0\ \exists\delta>0$,
	      使得
	      \begin{equation*}
	      \norm{\tau_hf-f}_p<\varepsilon\quad\forall f\in\mathcal{F}\
	      \text{且}\ \forall h\in\R,\ \abs{h}<\delta.
	      \end{equation*}
	\item 证明 $\mathcal{F}$ 在 $L^p(\R)$ 中不具有紧闭包.
\end{enumerate}
\end{exercise}

\begin{exercise}
\label{exer:4.34}
设 $1\le p<\infty$, $\mathcal{F}\subset L^p(\R^N)$ 为 $L^p(\R^N)$ 的紧子集.
\begin{enumerate}[label=\arabic*., leftmargin=2.2em, itemsep=0.35em]
	\item 证明 $\mathcal{F}$ 在 $L^p(\R^N)$ 中有界.
	\item 证明 $\forall\varepsilon>0\ \exists\delta>0$, 使得
	      \begin{equation*}
	      \norm{\tau_hf-f}_p<\varepsilon\quad\forall f\in\mathcal{F}\
	      \text{且}\ \forall h\in\R^N,\ \abs{h}<\delta.
	      \end{equation*}
	\item 证明 $\forall\varepsilon>0\ \exists\Omega\subset\R^N$ 有界, 开, 使得
	      \begin{equation*}
	      \norm{f}_{L^p(\R^N\setminus\Omega)}<\varepsilon\quad\forall f\in\mathcal{F}.
	      \end{equation*}
	      与推论  \ref{cor:4.27} 作比较.
\end{enumerate}
\end{exercise}

\begin{exercise}
\label{exer:4.35}
固定函数 $G\in L^p(\R^N)$, $1\le p<\infty$, 并令 $\mathcal{F}=G\star B$, 其中
$B$ 是 $L^1(\R^N)$ 中的有界集.

证明对任一测度有限的可测集 $\Omega\subset\R^N$, $\mathcal{F}|_\Omega$ 在
$L^p(\Omega)$ 中的闭包紧.与推论  \ref{cor:4.28} 作比较.
\end{exercise}

\begin{exercise}[等度可积族]
\label{exer:4.36}
若子集 $\mathcal{F}\subset L^1(\Omega)$ 满足下述性质, 则称它是\emph{等度可积}
的:\footnote{可以证明: 若测度空间 $\Omega$ 是弥散的 (即 $\Omega$ 没有原子), 则
(a) 由 (b) 与 (c) 推出.例如考虑带 Lebesgue 测度的 $\Omega=\R^N$.}
\begin{enumerate}[label=(\alph*), leftmargin=2.4em, itemsep=0.15em]
	\item $\mathcal{F}$ 在 $L^1(\Omega)$ 中有界,
	\item $\begin{cases}
	      \forall\varepsilon>0\ \exists\delta>0\ \text{使得}\\
	      \displaystyle\int_E\abs{f}<\varepsilon\quad\forall f\in\mathcal{F},\
	      \forall E\subset\Omega\ \text{可测, 且}\ \abs{E}<\delta,
	      \end{cases}$
	\item $\begin{cases}
	      \forall\varepsilon>0\ \exists\omega\subset\Omega\ \text{可测, 且}\
	      \abs{\omega}<\infty\ \text{使得}\\
	      \displaystyle\int_{\Omega\setminus\omega}\abs{f}<\varepsilon\quad
	      \forall f\in\mathcal{F}.
	      \end{cases}$
\end{enumerate}
设 $(\Omega_n)$ 为 $\Omega$ 中的非减可测集序列, $\abs{\Omega_n}<\infty\ \forall n$,
且 $\Omega=\bigcup_n\Omega_n$.
\begin{enumerate}[label=\arabic*., leftmargin=2.2em, itemsep=0.35em]
	\item 证明 $\mathcal{F}$ 等度可积当且仅当
	      \begin{equation*}
	      \lim_{t\to\infty}\sup_{f\in\mathcal{F}}\int_{[\abs{f}>t]}\abs{f}=0
	      \end{equation*}
	      且
	      \begin{equation*}
	      \lim_{n\to\infty}\sup_{f\in\mathcal{F}}\int_{\Omega\setminus\Omega_n}
	      \abs{f}=0.
	      \end{equation*}
	\item 证明: 若 $\mathcal{F}\subset L^1(\Omega)$ 紧, 则 $\mathcal{F}$ 等度
	      可积.逆命题是否成立?
\end{enumerate}
\end{exercise}

\begin{exercise}
\label{exer:4.37}
固定函数 $f\in L^1(\R)$, 使得
\begin{equation*}
\int_{-\infty}^{+\infty}f(t)\,dt=0\quad\text{且}\quad
\int_0^{+\infty}f(t)\,dt>0,
\end{equation*}
并对 $x\in I=(-1,+1)$ 令 $u_n(x)=nf(nx)$.
\begin{enumerate}[label=\arabic*., leftmargin=2.2em, itemsep=0.35em]
	\item 证明
	      \begin{equation*}
	      \lim_{n\to\infty}\int_Iu_n(x)\phi(x)\,dx=0\quad
	      \forall\phi\in C([-1,+1]).
	      \end{equation*}
	\item 验证序列 $(u_n)$ 在 $L^1(I)$ 中有界.证明 $(u_n)$ 的任一子列都不等度
	      可积.
	\item 证明不存在函数 $u\in L^1(I)$, 使得沿某个子列 $(u_{n_k})$ 有
	      \begin{equation*}
	      \lim_{k\to\infty}\int_Iu_{n_k}(x)\phi(x)\,dx
	      =\int_Iu(x)\phi(x)\,dx\quad\forall\phi\in L^\infty(I).
	      \end{equation*}
	\item 与 Dunford--Pettis 定理作比较 (见问题 23 的 A3 问).
	\item 证明存在子列 $(u_{n_k})$, 使得当 $k\to\infty$ 时 $u_{n_k}(x)\to0$
	      在 $I$ 上几乎处处成立.
	      [提示: 计算 $\int_{[n^{-1/2}<\abs{x}<1]}\abs{u_n(x)}\,dx$, 并应用定理
	       \ref{thm:4.9}.]
\end{enumerate}
\end{exercise}

\begin{exercise}
\label{exer:4.38}
令 $I=(0,1)$, 并考虑由
\begin{equation*}
u_n(x)=\begin{cases}
n & \text{若}\ x\in\displaystyle\bigcup_{j=0}^{n-1}
\Bigl(\frac jn,\frac jn+\frac1{n^2}\Bigr),\\[10pt]
0 & \text{其他情形}
\end{cases}
\end{equation*}
定义的 $L^1(I)$ 中的函数序列 $(u_n)$.
\begin{enumerate}[label=\arabic*., leftmargin=2.2em, itemsep=0.35em]
	\item 验证 $\abs{\operatorname{supp}u_n}=\frac1n$, $\norm{u_n}_1=1$.
	\item 证明
	      \begin{equation*}
	      \lim_{n\to+\infty}\int_Iu_n(x)\phi(x)\,dx=\int_I\phi(x)\,dx
	      \quad\forall\phi\in C([0,1]).
	      \end{equation*}
	      [提示: 先考虑 $\phi\in C^1([0,1])$ 的情形.]
	\item 证明 $(u_n)$ 的任一子列都不等度可积.
	\item 证明不存在函数 $u\in L^1(I)$, 使得沿某个子列 $(u_{n_k})$ 有
	      \begin{equation*}
	      \lim_{k\to\infty}\int_Iu_{n_k}(x)\phi(x)\,dx
	      =\int_Iu(x)\phi(x)\,dx\quad\forall\phi\in L^\infty(I).
	      \end{equation*}
	      [提示: 取进一步的子列 $(u_{n'_k})$, 使
	      $\sum_k\abs{\operatorname{supp}u_{n'_k}}<1$.]
	\item 证明存在子列 $(u_{n_k})$, 使得当 $k\to\infty$ 时 $u_{n_k}(x)\to0$
	      在 $I$ 上几乎处处成立.
\end{enumerate}
\end{exercise}
